\documentclass[11pt,a4paper]{article}
\usepackage[T1]{fontenc}
\usepackage[utf8]{inputenc}
\usepackage{mathptmx}
\usepackage{amsmath}
\usepackage{graphicx}
\usepackage{booktabs}
\usepackage{array}
\usepackage{adjustbox}
\usepackage{placeins}
\usepackage[margin=2.5cm]{geometry}
\usepackage[authoryear,round]{natbib}
\usepackage{microtype}
\usepackage[font=small,labelfont=bf]{caption}
\usepackage{xurl}
\usepackage[hidelinks]{hyperref}
\usepackage{authblk}
\urlstyle{same}
\usepackage{lineno}
\linenumbers
\renewcommand{\linenumberfont}{\normalfont\tiny\color{gray}}
\usepackage{xcolor}

\title{The cost-optimal share of nuclear power in a zero-emission Italian electricity system: a capacity expansion analysis at three-hour resolution with auction prices}
\author{Claudio Gratton}
\affil{Independent researcher, Italy}
\date{Preprint, November 2026}

\begin{document}
\maketitle

\begin{abstract}
\noindent Italy approved in 2026 a law to reintroduce nuclear power, while the public debate still compares technologies through their levelised cost of electricity. We estimate the least-cost zero-emission Italian power system for a 2040 demand of 400\,TWh with a capacity expansion model in which one fleet serves four weather years and every technology that can produce, store or shift electricity is a decision variable under one cost rule; technologies without commercial plants have a first-plant (PI) and a target (OB) case. Imports follow a rule of net self-sufficiency: in each year Italy imports no more than it exports, and it gets no help from abroad in the three critical weeks of each year, so that the plan would still work if every country planned the same way. The result is one curve: the annualised fixed cost of nuclear (fixN, EUR/kW-yr, combining construction cost, cost of capital and construction time) against optimal nuclear capacity. Cost cases are points on it. With the Dukovany cost (contract price net of fuel plus owner's costs, 8{,}540\,EUR/kW, 5\% cost of capital, seven years, fixN 670), the three-hour central system has <<RUN:ss0_PI_670_h3:N_GW|1>>\,GW of nuclear (<<RUN:ss0_PI_670_h3:sh|1>>\% of demand) in PI and <<RUN:ss0_OB_670_h3:N_GW|1>>\,GW (<<RUN:ss0_OB_670_h3:sh|1>>\%) in OB, and the gap between the system without and with nuclear is <<GAP:ss0_PI_none_h3:ss0_PI_670_h3|2>> and <<GAP:ss0_OB_none_h3:ss0_OB_670_h3|2>>\,EUR/MWh. Allowing net imports of up to 10\% of demand saves only <<EXPR:V('RUN:ss0_PI_670_h6:cost_new')-V('RUN:ss10_PI_670_h6:cost_new')|2>> (PI) and <<EXPR:V('RUN:ss0_OB_670_h6:cost_new')-V('RUN:ss10_OB_670_h6:cost_new')|2>>\,EUR/MWh (OB) with nuclear, against <<EXPR:V('RUN:ss0_PI_none_h6:cost_new')-V('RUN:ss10_PI_none_h6:cost_new')|2>> and <<EXPR:V('RUN:ss0_OB_none_h6:cost_new')-V('RUN:ss10_OB_none_h6:cost_new')|2>> without it. The six-hour step underestimates the value of nuclear, so headline values are three-hour results; a Richardson estimate puts the PI gap near <<RICHS:PI:670:G1|1>>\,EUR/MWh towards a one-hour step, while for OB the three-hour gap is a lower bound. At six hours nuclear falls below 10\,GW above a fixN of about <<THR6:PI:s:5:10|-1>> (PI) and <<THR6:OB:s:5:10|-1>> (OB), about <<EXPR:V('CPX:'+str(V('THR6:PI:s:5:10'))+':5:7')|-2>> and <<EXPR:V('CPX:'+str(V('THR6:OB:s:5:10'))+':5:7')|-2>>\,EUR/kW with seven years of construction. Two conditions halve the result or more, both through the fixed cost: at a 7\% cost of capital the central case gives <<RUN:ss0_PI_7pc_923_h3:N_GW|0>>\,GW in PI and <<RUN:ss0_OB_7pc_923_h3:N_GW|0>> in OB at three hours. Fifteen-year construction at the central cost gives <<N6:PI:s:5:792.5|0>> and <<N6:OB:s:5:792.5|0>>\,GW at six hours; worst weather years raise the optimum to about <<RUN:ss0_wx4_PI_670_worst:N_GW|0>>\,GW. Under the earlier rule, imports with a capacity credit and no exports, the central case had <<RUN:c3_PI_670_h3:N_GW|0>>\,GW (PI). About <<EXPR:V('RUN:ss0_PI_670_h3:N_GW')/1.055|0>> APR1000 units are needed; the construction pace is a constraint to verify.
\end{abstract}

\noindent\textbf{Keywords:} nuclear power; capacity expansion; Italy; system cost; cost of capital; import dependence; variable renewables

\medskip
\noindent\textbf{Highlights}
\begin{itemize}\setlength\itemsep{0pt}
\item One curve links the annualised fixed cost of nuclear to optimal capacity; cost cases are points on it.
\item Imports follow net self-sufficiency: no net imports over a year and no foreign help in critical weeks.
\item Central case (8{,}540\,EUR/kW including owner's costs): <<RUN:ss0_PI_670_h3:N_GW|1>>\,GW (<<RUN:ss0_PI_670_h3:sh|1>>\% of demand) in PI and <<RUN:ss0_OB_670_h3:N_GW|1>>\,GW (<<RUN:ss0_OB_670_h3:sh|1>>\%) in OB at three hours.
\item With nuclear, a 10\% net import allowance saves about 1--1.5\% of system cost: a poor target.
\item A 7\% cost of capital leaves <<RUN:ss0_PI_7pc_923_h3:N_GW|0>> and <<RUN:ss0_OB_7pc_923_h3:N_GW|0>>\,GW; the six-hour step underestimates nuclear.
\end{itemize}

\section{Introduction}\label{sec:intro}

In September 2026 the Italian Parliament gave final approval to a law delegating to the government the reintroduction of nuclear power \citep{Delega2026}. The law does not authorise any plant and does not select sites; the national energy and climate plan foresees 8--16\,GW of nuclear capacity by 2050, equal to 11--22\% of electricity demand. How much nuclear capacity a zero-emission Italian power system should contain, and under which conditions, is therefore a practical question.

The public debate in Italy relies mostly on the levelised cost of electricity (LCOE). In the IEA World Energy Outlook 2024, utility-scale solar in the European Union in 2030 costs about 35\,USD/MWh against about 140\,USD/MWh for nuclear \citep{IEA2024WEO}; the same comparison is common in the Italian debate \citep{Pagella2026}. Industry reports such as \citet{Lazard2025} and \citet{Fraunhofer2024} reach similar conclusions. The LCOE, however, divides the cost of a plant by the energy it produces regardless of when it produces it. In Italy solar output in January is one third of July output (Section~\ref{sec:monthly}), so a solar MWh produced at noon in July and a nuclear MWh produced on a January evening are not the same product. System studies address this. \citet{Sepulveda2018}, analysing nearly 1{,}000 cases for US systems, find that firm low-carbon resources reduce the cost of electricity by 10--62\% in fully decarbonised cases, which is an increase of 11--163\% in the cost without them (our conversion); the \citet{NEA2019} estimates integration costs of variable renewables rising from about 8 to about 50\,USD/MWh between 10\% and 75\% penetration.

For Italy, \citet{Lombardi2020} find a near-optimal 2050 system with 144\,GW of solar, 76\,GW of wind and synthetic methane for back-up, without nuclear. \citet{Mastrantuono2026} find that new nuclear lowers total system cost by up to 8 billion euros, roughly halves imports and reduces storage needs by 42\%, curtailment by 35\% and land use by 25\%; their model, however, caps nuclear at 8\,GWe, and the cap binds in the scenarios with stable costs, so the optimal nuclear share is not identified. An analysis by Politecnico di Milano and Universit\`a dell'Insubria, reported in the trade press, finds up to 15\,GW of nuclear at 5{,}500\,EUR/kW and none at 8{,}300\,EUR/kW with ten-year construction; the model has not been published \citep{PoliMi2026}.

This paper contributes in five ways. First, it estimates the cost-optimal amount of nuclear in a zero-emission Italian system without exogenous caps, for a 2040 demand that includes heat pumps and electric vehicles, with a model at three-hour resolution in which one fleet must serve four weather years, as a planner would require, and with an explicit adequacy reserve. Second, it prices renewables and storage from the prices actually awarded in Italian auctions, rather than from engineering estimates, and converts these nominal, partly indexed contract prices into real costs over the life of the plants, so that they are comparable with a nuclear cost computed in real terms over its own life. Solar and wind follow finite supply curves, because ground-mounted modules are banned on agricultural land and wind sites are limited, so that land enters the cost of the system; existing plants enter at operating cost only. Third, it makes the dependence of the result on the three variables that drive nuclear cost, construction cost, cost of capital and construction time, explicit by summarising them in a single quantity, the annualised fixed cost per kW (fixN), so that the result is one curve and each cost case is a point on it. Fourth, it includes every technology that can produce, store or shift electricity in Italy under one cost rule, with first-plant and target cases for technologies without commercial plants. Fifth, it treats imports as a dependence choice with a price: the central rule is net self-sufficiency, with exports allowed and no foreign help in the critical weeks, and the cost saved by allowing net imports is reported separately, so that a plan is not made to look cheaper by counting on plants that other countries build. The model has one node and no ramp constraints: the aim is to identify the variables that matter and to make the result reproducible, not to represent every feature of the Italian system. It describes a regime system, with 2040 demand and all fleets already built, not the situation in 2040. The result also supersedes an earlier version of this analysis that gave 38.5\,GW and 59\% of demand, in a model without imports, hydrogen, bioenergy, geothermal and the other technologies now included; Section~\ref{sec:bridge} explains the difference step by step. Section~\ref{sec:methods} describes the model and the data, Section~\ref{sec:results} presents the results, Section~\ref{sec:discussion} discusses them together with supply security, industrial capability and political feasibility, and Section~\ref{sec:conclusions} concludes.

\section{Methods and data}\label{sec:methods}

\subsection{Metric}
We compare options by the cost of the electricity that actually meets demand in every period, rather than by the levelised cost of the energy each plant produces. For a single technology on a monthly basis this is the levelised cost multiplied by the overbuild needed to cover the worst month, plus the cost of the storage that shifts energy into the hours and months in which it is missing (Section~\ref{sec:monthly}). For the system as a whole it is the total annualised cost of the least-cost portfolio that meets demand in every time step, divided by the demand not covered by hydropower, \texttt{cost\_new} in the repository. This single denominator (356.6\,TWh of the 400\,TWh) is used throughout the system results; multiplied by 356.6/400 = 0.8915 it gives the cost per MWh of total demand, for example <<EXPR:V('RUN:ss0_PI_670_h3:cost_new')*0.8915|1>> instead of <<RUN:ss0_PI_670_h3:cost_new|2>>\,EUR/MWh for the central PI case at three hours. The monthly calculation divides instead by total demand. Monthly and hourly costs must be compared only through this conversion. Because the time step biases absolute costs (Section~\ref{sec:step}), the decision metric is the gap, the cost of the system without nuclear minus the cost of the system with nuclear.

\subsection{Capacity expansion model}\label{sec:model}
The model is a linear programme that chooses the capacity of nuclear ($N$), of each block $b$ of the solar ($S_b$), onshore wind ($W_b$) and offshore wind ($O_b$) supply curves, and the energy capacity of lithium-ion batteries ($E$), and their dispatch, so as to meet demand in every time step at least cost, with zero direct emissions (mode ZD). Version 7 adds every other technology that can produce, store, shift or reduce electricity in Italy as a decision variable with one uniform cost rule: imports (11.5\,GW existing and 7.3\,GW new), existing and new hydropower and pumped hydro, geothermal (existing, two conventional blocks and enhanced systems), bioenergy (biogas, solid biomass, waste), concentrated solar power with storage, wave energy, compressed air, hydrogen (electrolysis, geological storage, turbine), 100-hour storage, smart charging, interruptible load and small modular reactors; gas with carbon capture and direct air capture enter only in the net-zero (NZ) variants. No technology is excluded without a number showing that it is negligible (tidal currents, for example, are 0.031\% of demand). Technologies with market prices or contracts use the latest observed price. Technologies without commercial plants have two cases, both reported: PI (first plant, costs of demonstrators or first offers) and OB (target, costs declared by developers or projected). The OB model is reduced: it has no wave energy, 100-hour storage, new pumped hydro, new hydropower or SMR. The time step is six hours for the curve points and three hours for the corrected central values (Section~\ref{sec:step}); the 3-hour data are block sums, so a one-hour step is not possible. The four weather years 2016--2019 are optimised jointly: one set of capacities must serve all of them, with dispatch optimised in each year. A power system is built before the weather is known, and optimising each year on its own would give every year a fleet designed for its own weather. The model has a single node, so energy moves between regions without losses or constraints beyond a grid charge on capacity, and it has no ramp constraints: nuclear output may move between its minimum and its maximum from one step to the next at no cost. The objective is
\begin{equation}
\min \; Y\Big(c_N N + \sum_b c_{S_b} S_b + \sum_b c_{W_b} W_b + \sum_b c_{O_b} O_b + c_E E\Big) + \sum_{t} v_N\, n_t ,
\end{equation}
where $Y$ is the number of years optimised together and $c_i$ the annualised cost of capacity $i$, including the grid charge. For solar and wind the cost of a block is the real lifetime cost derived from auction prices (Section~\ref{sec:inputs}) multiplied by the full-load hours of the block, so that each MWh the block can produce is paid whether or not it is used. Existing solar plants installed after 2016 and the existing wind block E0 enter at operating cost only. The energy balance in each step $t$ of length $\Delta$ is
\begin{equation}
n_t + p^S_t + p^W_t + p^O_t + h_t + d^B_t - q^B_t + d^P_t - q^P_t = L_t ,
\end{equation}
where all flows are energies in step $t$, $L_t$ is the demand net of the flat generation from bioenergy and geothermal, $n_t$ nuclear output, $p^S_t \le \sum_b S_b\,\rho_b\,s_t$ the solar output used ($s_t$ the PECD energy per MW and $\rho_b$ the scale of the block for its zone), with the same form for onshore and offshore wind, so that curtailment is allowed; $h_t$ is hydropower and $d$, $q$ are discharge and charge of batteries ($B$) and pumped hydro ($P$). Nuclear output follows load between a minimum and a maximum share of the available capacity,
\begin{equation}
\mu\,a_t\,N\Delta \le n_t \le a_t\,N\Delta ,
\end{equation}
with minimum load $\mu=0.5$ and availability $a_t$ that depends on the month: 0.97 in December, January, February, June, July and August and 0.87 in the other six months, a mean of 0.92. This availability is an assumption: it is built to match the 92\% load factor of the US fleet \citep{EIA2025} and refuelling outages concentrated in spring and autumn \citep{EIA2023}, not a measurement of an Italian fleet. Storage follows
\begin{equation}
e^B_{t+1} = e^B_t + \eta_B q^B_t - d^B_t,\quad 0\le e^B_t \le E,\quad q^B_t, d^B_t \le E\Delta/\tau ,
\end{equation}
with round-trip efficiency $\eta_B=0.85$ and duration $\tau=6.6$\,h, and the same form for pumped hydro with fixed power 7.3\,GW, energy 53\,GWh and $\eta_P=0.75$; the other storage technologies and imports enter the balance with their own flows and costs (Table~\ref{tab:inputs}). Storage levels are cyclic. Hydropower in each month $m$ must produce the monthly energy $H_m$, with a minimum equal to a share $f=0.37$ of the monthly mean (run-of-river flow \citep{TernaHydro}) and a maximum power $\bar h=15$\,GW; the budgets $H_m$ are the same in the four simulated years, so hydropower is independent of the simulated wind and solar weather:
\begin{equation}
\sum_{t\in m} h_t = H_m,\qquad f\,\frac{H_m}{|m|}\le h_t \le \bar{h}\Delta .
\end{equation}

\paragraph{Adequacy reserve.} A static constraint requires the credited capacity to exceed the net peak $\hat L$ (66.98\,GW at six hours, 76.88\,GW at three hours) by a factor $\rho=1.1$:
\begin{equation}
a_{\mathrm{pk}} N + E/\tau + \bar h_c + \bar P + \kappa I \ge \rho \hat L ,
\end{equation}
where $a_{\mathrm{pk}}=0.97$ is the nuclear availability in the month of the peak, $\bar h_c=15$\,GW the hydro credit, $\bar P=7.3$\,GW the pumped hydro, $I$ the import capacity and $\kappa$ its credit: zero in the central case, 0.38 in the variant of the earlier version (4.4 of 11.5\,GW in the 2025 capacity market); credits of the other technologies are in the repository. All credits are declared assumptions.

\paragraph{Imports and net self-sufficiency.} Imports $i_t$ use the interconnection capacity $I$ (11.5\,GW existing and up to 7.3\,GW new) at 115.94\,EUR/MWh, the 2025 average PUN \citep{GME2025}. Exports $x_t\le I\Delta$ use the same capacity and earn nothing, a cautious assumption for the system that exports, because surplus energy is worth little precisely when it is also abundant in neighbouring countries. The central rule is net self-sufficiency: for every weather year $y$,
\begin{equation}
\sum_{t\in y} i_t \le \sum_{t\in y} x_t + X\,D_y ,\qquad X=0 ,
\end{equation}
where $D_y=400$\,TWh, and in the three critical weeks of each year imports are not available and $\kappa=0$. The rule asks that the Italian plan would still work if every European country planned the same way: an Italy that covers a fixed share of demand with energy produced abroad leans on plants that others build and pay for, much of them French nuclear. Variants: $X=5\%$ and $10\%$; imports without an annual bound and without exports but with no help in the critical weeks; and the rule of the earlier version, imports without exports, a 17\% price mark-up in the critical weeks and an availability and credit of 38\%. With $X=100\%$ and the 38\% credit the model reproduces the earlier central case to the cent. The model does not simulate the neighbouring systems. The twenty highest net-load steps all fall in winter evenings between 17:00 and 21:00, mostly in January, which is why heat pumps and electric vehicles push towards firm capacity.

\paragraph{Nuclear cost.} The annualised fixed cost combines overnight construction cost $K$, real cost of capital $r$ and construction time $T_c$:
\begin{equation}
c_N = K\,\phi(r,T_c)\,\mathrm{CRF}(r,L_N) + F_N ,\qquad
\phi(r,T_c)=\frac{1}{T_c}\sum_{k=0}^{T_c-1}(1+r)^{T_c-k-1/2},\qquad \mathrm{CRF}(r,n)=\frac{r(1+r)^n}{(1+r)^n-1},
\end{equation}
where $\phi$ accounts for interest during construction with uniform spending, $L_N=80$ years, and $F_N$ is the fixed part: operation at 17\,EUR/MWh at 90\% load, $17\times0.9\times8760=134.03$\,EUR/kW-yr, plus a grid charge of 2\,EUR/MWh, 16.11\,EUR/kW-yr, so $F_N=150.14$ at 5\% and 154.77 at 7\%, because at 7\% the grid charge is multiplied by 1.287. We call $c_N$ \emph{fixN}. The inverse is $K=(\mathrm{fixN}-F_N)/(\phi\,\mathrm{CRF})$. The variable cost is 15\,EUR/MWh. At 7\% the capital part of every technology priced by auction or by capex is multiplied by $\mathrm{CRF}(7\%,L)/\mathrm{CRF}(5\%,L)$ for its life $L$ (1.176 for 20 years, 1.239 for 30, 1.287 for 40); operation, fuel and the price of imported energy are unchanged, and nuclear uses its own $r$. The model is solved with HiGHS through its Python interface; all solutions have status optimal. The validation set is five reference runs (no-nuclear PI and OB, PI and OB at fixN 598 with imports, PI at 817 with 7\%), reproduced within 0.05\,EUR/MWh and 0.2\,GW.

\subsection{Data and assumptions}\label{sec:inputs}
Table~\ref{tab:inputs} lists the inputs with their status: DATA (read in a source), DERIVED (calculation on data), ESTIMATE (the author's judgement with a declared method) or ASSUMPTION (a model choice without a source, with a sensitivity).

\paragraph{Demand.} The 2040 demand is 400\,TWh without electrolysers. It is the sum of the ENTSO-E hourly load of 2016--2019 \citep{OPSD2020}, rescaled to 338\,TWh in each year, 15\,TWh of heat pumps and 47\,TWh of electric vehicles. The value is a central choice between the 385\,TWh for 2035 and the 415 and 439\,TWh for 2040 of the Terna and Snam scenario report \citep{DDS2024}. Heat pumps follow Italian monthly heating degree-days (Eurostat, \texttt{nrg\_chdd\_m}); their hourly shape, with morning and evening peaks, is an assumption. Electric vehicles charge 70\% of their daily energy between 18:00 and 24:00 and 30\% spread over the day, an assumption tested with flat, night-only, evening-only, daylight (50\% in daylight hours) and mixed profiles. Bioenergy (8\,TWh net) and geothermal (5.3\,TWh net) are subtracted as flat generation.

\paragraph{Weather.} Solar, onshore and offshore wind capacity factors are hourly PECD 2021.3 profiles for the 2030 fleet in Italy \citep{PECD2021}, used without normalisation: the mean full-load hours of 2016--2019 are 1{,}375 for solar, 2{,}227 for onshore wind and 2{,}935 for offshore wind, and differences between years are kept. For the worst-weather check we use the 38 PECD years 1982--2019 (Section~\ref{sec:results}).

\paragraph{Hydropower.} Annual hydro energy is 43.4\,TWh, the 2016--2019 mean excluding pumping, with the monthly shape of the 2023--2024 mean \citep{Terna2025m}. Pumped hydro power is 7.3\,GW (Terna 2024) and its energy 53\,GWh (Ambrosetti-Edison); efficiency is an assumption.

\paragraph{Solar supply curve.} Solar has 13 blocks by type (ground-mounted, industrial rooftops, other rooftops, agrivoltaic) and zone (South and islands, Centre-South, North), with a limit in GW, a scale $\rho_b$ for the yield of the zone from PVGIS 5.3 \citep{PVGIS}, and a cost multiplier on the capital part. Ground-mounted modules are banned on agricultural land by Article 5 of Decree-Law 63/2024 \citep{DL63}, a rule upheld by Constitutional Court judgment 127/2026 and taken up by Article 11-bis introduced by Decree-Law 175/2025 \citep{Consulta127}. No national estimate of the remaining non-agricultural area exists, so the limits are our estimates anchored to ISPRA \citep{ISPRA2026}: in the central curve 42.7\,GW on the ground, 140\,GW of rooftops (the lower end of the 140--193\,GW range of ISPRA 2026, of which 30\,GW industrial at 1{,}000\,EUR/kW) and 119.7\,GW of agrivoltaic (3\% of suitable agricultural area at 0.5\,MW per hectare), 302\,GW in all. For comparison \citet{Ranjgar2025} estimate 28--52\,GW of ground-mounted potential with 1\% of the available area. Ground-mounted plants occupy 1{,}700\,ha/GW \citep{ISPRA2026}. Output per kW for new plants is 10--16\% higher in the South than in the North (PVGIS, Sicily to Lombardy 1.16, Apulia to Veneto 1.10); the larger difference in the yield of the existing fleet mixes orientation, age and curtailment \citep{GSE2025}. The installed fleet was 43.5\,GW at the end of 2025 and 19.3\,GW at the end of 2016 \citep{ItaliaSolareQ4}; the 24.2\,GW installed after 2016 enter at operating cost and are subtracted from the blocks. The capital share of the real cost of solar is 76.4\%; multipliers apply to that part: 1.333 for agrivoltaic (1.0 against 0.75 million EUR/MW, midpoint of the GSE range reported by \citet{QualEnergiaAgri}) and 1.5625 for rooftops other than industrial, within the 900--1{,}600\,EUR/kWp range of \citet{FraunhoferEN2024}.

\paragraph{Wind supply curve.} Onshore wind has six blocks: E0, 13.6\,GW of existing repowered plants at operating cost (31\,EUR/MWh, 1{,}600 h), E1 6\,GW at 64\,EUR/MWh (2{,}400 h), E2 10\,GW at 71 (2{,}100 h), E3 15\,GW at 82 (1{,}800 h), E4 15\,GW at 96 (1{,}500 h) and E5 10\,GW at 118 (1{,}200 h), 69.6\,GW in all; block scale $\rho_b$ is hours of the block over the 2{,}227 h of the PECD profile. Costs follow JRC-EU-TIMES parameters (1{,}320\,kEUR/MW, 5.2\%, 25 years) \citep{JRCEUTIMES}, checked against the FER~X auction at 72.85\,EUR/MWh \citep{FERXmain2025}. The potential for Italy in the same source ranges from 10.1 to 171.6\,GW with the planning rules used. Offshore wind has five blocks, O1 2\,GW at 160\,EUR/MWh, O2 3.8\,GW at 172.5, O3 6\,GW at 138.5, O4 10\,GW at 112.5 and O5 25\,GW at 132.5, 46.8\,GW in all, with grid connection and export cable included in the cost and so no base grid charge.

\paragraph{Renewable and storage prices.} The first FER~X auction (December 2025) awarded 7.7\,GW of solar at an average of 56.82\,EUR/MWh \citep{FERXmain2025}. For batteries we use the first MACSE auction of 2025, 12{,}959\,EUR per MWh of capacity per year for 6.6-hour lithium-ion batteries and 15-year contracts \citep{MACSE2025}. The MACSE premium is the fixed remuneration paid to the holder of the capacity: market revenue from time shifting goes to the operator Terna and reduces the burden on consumers, and the holder keeps only a share of the margins from ancillary services \citep{ARERA247}; the premium is therefore a good estimate of the full capacity cost, if anything slightly low. FER~X prices are indexed to inflation for 20\% of the strike price after commissioning (art.\ 11 c.\ 3 of DM 457/2024) \citep{FERXrules2025,FERXindex}; the 20\% applied to every technology is an assumption, and the MACSE premium is indexed for 10\% \citep{MACSErules}. We convert each auction price $p$ into a real levelised cost over the plant life,
\begin{equation}
c = p\,\frac{\sum_{t=1}^{T}\left[\iota+(1-\iota)(1+\pi)^{-t}\right](1+r)^{-t}}{\sum_{t=1}^{T}(1+r)^{-t}}\left[(1-o)\frac{\mathrm{CRF}(r,L)}{\mathrm{CRF}(r,T)}+o\right],\label{eq:conv}
\end{equation}
with contract length $T$, indexed share $\iota$, inflation $\pi=2\%$, real discount rate $r=5\%$ for renewables and storage (also in the cases where nuclear is financed at 4\%), operating-cost share $o$ and plant life $L$. The real cost of solar is 42.2\,EUR/MWh, composed of capital (76.4\%) and operation (9.94\,EUR/MWh), and of batteries 11{,}466\,EUR/MWh-yr; indexation of 0\% or 100\% would give 40.7 or 48.2\,EUR/MWh (our calculation). Disposal is assumed to be in the offer price, because the rules put it on the producer or operator: the electronic-waste rules for solar \citep{RAEE}, the restoration guarantee required by the permit for wind \citep{WindRestoration} and extended producer responsibility for batteries \citep{BatteryReg}; the estimated effect is 1--2\,EUR/MWh. Plants are rebuilt every 30 years (solar and wind) and 15 years (batteries). Existing plants after 2016 and E0 are assumed not to need renewal before 2040.

\paragraph{Grid.} The grid charge is on new capacity: 13\,EUR/MWh for solar and 8 for onshore wind (a charge per MW-year equal to those values times 1{,}066 and 1{,}709 h) and 2\,EUR/MWh for nuclear, doubling on new solar and wind above 137\,GW, the capacity enabled by the 2035 transmission plan. The values derive from Terna's development plan, 23 billion euros for 94\,GW of new capacity, about 245\,EUR/kW, with a 5\% annuity over 40 years, our rounding of the ARERA conventional lives of 45 years for transmission lines and 33 for stations \citep{ARERA615}, applied to the plan of \citet{TernaPdS2025}; the doubling follows the increasing integration cost of \citet{Ueckerdt2013}, and the factor 2 is an estimate. At 7\% the grid charge is multiplied by 1.287 (our assumption for the scale of the grid).

\paragraph{Nuclear: cost cases.} Cost cases are points on one curve of fixN, not separate runs: Saeul 2{,}965\,EUR/kW with 8 years (fixN 335.2 at 5\%), Barakah 4{,}670 with 8 years (<<FIX:4670:5:8|1>>), the central case Dukovany 8{,}540 with 7 years (669.5 at 5\%, 922.7 at 7\%; the runs are at 670 and 923), a risk case 9{,}680 with 15 years (878.3 at 5\%, 1{,}334.2 at 7\%) and Vogtle 15{,}540 with 10 years (1{,}172.2). The central case stands for plants built properly, with one series design and engineering completed before construction starts; Flamanville, Olkiluoto 3, Mochovce and Vogtle are the Western first-of-a-kind builds behind the risk case, and Hinkley Point C is another risk reference that is not in its median. All five cases are overnight construction costs, without interest during construction (which the model adds through the factor idc), in 2025 euros per kW and including owner's costs. Dollar values are brought to 2025 with the US consumer price index (CPI-U, 274.310 in September 2021 and 321.943 on average in 2025, eleven months, October is missing because of the federal funding lapse \citep{BLSCPI}) and converted at 1 EUR = 1.1188 USD, the Bank of Italy reference rate of 26 October 2026 \citep{BancaItaliaFX}. The annual CPI-U averages before 2016, needed for Barakah and Vogtle, are the means of the twelve monthly values of the not seasonally adjusted series (BLS CUUR0000SA0, retrieved from FRED \citep{FREDCPI}): 214.537 in 2009, 218.056 in 2010, 224.939 in 2011, 229.594 in 2012, 232.957 in 2013, 236.736 in 2014 and 237.017 in 2015. The mean factors for Barakah and Vogtle recomputed with them are 1.2549 and 1.2623, equal to the values used (1.255 and 1.262), so no value changes.

\paragraph{Nuclear: Saeul and Barakah.} Saeul is the Korean APR1400 series (about 3{,}300\,USD/kW for a 1.4\,GW unit \citep{HSF2026,KHNP2025}). Barakah is the same design exported to the United Arab Emirates: the construction contract is about 20.4 billion USD for four units, 5.6\,GW \citep{WNAuae}, that is 3{,}643\,USD/kW nominal. The 24.4 billion USD of total financing in the same source includes interest and non-construction items and is not used, because the model adds interest itself. With uniform spending over 2012--2024 (mean CPI-U factor 1.255) the contract is 4{,}571\,USD/kW in 2025 dollars, 4{,}086\,EUR/kW. As at Dukovany the contract is a construction price (it also includes the first fuel loads, which we do not remove) and does not include owner's costs; adding them with the rule below gives 4{,}673, stated as 4{,}670\,EUR/kW (our calculation), 7\% above the 4{,}371 that the total financing would give.

\paragraph{Nuclear: the central case, Dukovany.} The central value is built in four steps. First, the contract price: the Czech Republic signed with KHNP on 4 June 2025 the contract for two APR1000 units of 1{,}055\,MWe \citep{Dukovany2024,WNNDukovanySigned2025,EnerdataDukovany2025}, CZK\,407 billion, USD\,18.6 billion, for construction, commissioning and fuel for about ten years (first load and five reloads) according to World Nuclear News \citep{WNNKHNPplans}. The exchange rate is the 2025 annual average of the ECB reference rate, 24.688 CZK per EUR \citep{EurostatFX}; an earlier version assumed 25, and the correction raises the cost by about 1.3\%. The price is then 16.5 billion euros, 7{,}813\,EUR/kW for 2 $\times$ 1{,}055\,MWe and 7{,}850 for 2{,}100\,MW; the model uses 7{,}850. The price is a June 2025 price, so it is in 2025 prices and is not escalated; the indexation in the contract concerns only the nominal price. Second, the fuel. The sources disagree. World Nuclear News says the CZK\,407 billion include fuel for about ten years (first load and five reloads) \citep{WNNKHNPplans}, while \v{C}T24, quoting the director of Elektr\'arna Dukovany II, and Forbes \v{C}esko report CZK\,15 billion more for the first fuel \citep{CT24Dukovany,ForbesDukovany}. We adopt the first because it states what the price covers (first load and five reloads), whereas the other two do not say whether the CZK\,15 billion are inside or outside the CZK\,407 billion. The central case therefore removes the fuel, because the model already counts fuel in the 32\,EUR/MWh of operating costs: fuel costs 4.85\,EUR/MWh (uranium, conversion, enrichment and fabrication), times 10 years, 0.9 load and 8{,}760 hours divided by 1{,}000, that is 382\,EUR/kW, and 7{,}850 minus 382 is 7{,}468 (our calculation). The fuel cost is 4.62\,USD/MWh in September 2021 prices \citep{WNAecon}; 4.62 times the CPI-U factor 1.174 divided by 1.1188 is 4.85\,EUR/MWh. The runs use 4.91, a negligible difference (0.4\,EUR/kW-yr in fixN). Effect of the other reading: if the CZK\,407 billion did not include the fuel and it were not removed, the cost before owner's costs would be 7{,}850 instead of 7{,}468 (382\,EUR/kW more) and the overnight cost 8{,}980 instead of 8{,}540 (440\,EUR/kW, 5\% more; fixN 696 instead of 670).

Third, owner's costs. Neither the contract nor the sources that report it \citep{WNNKHNPplans,CT24Dukovany,ForbesDukovany} give them, and we found no public EDU~II estimate that separates them, so the contract price is an EPC price and the overnight cost is higher. We add them with the share that the Idaho National Laboratory reports for a pressurised-water reactor in the EEDB cost database: owner's costs (account 40) are 12.57\% of overnight cost \citep{INLowner}, so overnight cost is the EPC price divided by $(1-0.1257)$, that is multiplied by 1.1438, which is 14.4\% of the EPC price. As a European reference, the synthesis by D'haeseleer for the European Commission (DG Energy, 2013) notes that there is no unique definition and reports from several sources owner's costs of about 15--20\% of the EPC cost \citep{DHaeseleer2013}: our 14.4\% lies just below the low end of that range, and its high end of 20\% gives 8{,}962\,EUR/kW (7{,}468 $\times$ 1.20), fixN 695. The ends of the range are read off the curve and are not rerun; the central value is the INL share. Fourth, the result: $7{,}468\times1.1438=8{,}542$, stated as 8{,}540\,EUR/kW (our calculation), with fixN 669.5 at 5\% and 922.7 at 7\%.

The central value has three bounds, all on the same curve. The lowest cost is the EPC price net of fuel alone, 7{,}470\,EUR/kW, fixN 604: optimal nuclear is about <<N6:PI:s:5:604|1>>\,GW (PI, six hours) against <<RUN:ss0_PI_670_h6:N_GW|1>> in the central case at the same step; it is the upper bound for nuclear because it ignores costs that exist. The highest cost adds the state preparation spending of CZK\,85 billion quoted by \v{C}T24 \citep{CT24Dukovany}, 21\% of the CZK\,407 billion, instead of the 14.4\% mark-up: 9{,}030\,EUR/kW, fixN 699 (close to the 20\% of D'haeseleer), with <<N6:PI:s:5:699|1>>\,GW in PI and <<N6:OB:s:5:699|1>> in OB at six hours, against <<N6:PI:s:5:670|1>> and <<N6:OB:s:5:670|1>> in the central case at six hours. This figure is an assumption: we do not know how much of the CZK\,85 billion is construction. Without removing the fuel (8{,}980\,EUR/kW, fixN 696) the six-hour result is <<N6:PI:s:5:696|1>>\,GW in PI and <<N6:OB:s:5:696|1>> in OB. Owner's costs are the least documented item of the paper, and the choice of the central case depends on them.

\paragraph{Nuclear: risk case and Vogtle.} The risk case, 9{,}680\,EUR/kW with 15 years, is the median of four recent Western builds in 2025 euros per kW: Mochovce 3--4 about 8{,}760, Olkiluoto 3 about 9{,}060, Flamanville 3 about 10{,}300 and Vogtle 3--4 15{,}540; the median is the mean of the two central values, (9{,}060 + 10{,}300)/2 = 9{,}680 (our calculation). Olkiluoto 3 cost about 11 billion \citep{NEIOL3} for 1{,}600\,MWe (1{,}575 net \citep{WNAfinland}), 6{,}875\,EUR/kW nominal. The source reads ``\$11bn (\$12bn)'': we read 11 billion euros, 12 billion dollars; if it were 11 billion dollars, Olkiluoto 3 would be about 8{,}100\,EUR/kW and the median about 9{,}530 (1.6\% lower). Mochovce 3--4 cost 6.7 billion euros \citep{WNNMochovce4} for 2 $\times$ 471\,MWe, 7{,}113\,EUR/kW nominal (7{,}614 with the 440\,MWe net per unit of \citealt{WNAslovakia}). Brought to 2025 euros with the HICP \citep{EurostatHICP}, assuming uniform spending in 2005--2022 for Olkiluoto and 2009--2025 for Mochovce (mean factors 1.318 and 1.231), they are about 9{,}060 and 8{,}760\,EUR/kW. The sources give neither the spending schedule nor the price year: with other schedules Olkiluoto ranges from 8{,}907 to 9{,}208\,EUR/kW. Flamanville cost 13.2 billion euros of 2015, excluding interim interest \citep{EDFFlamanville2022}, for 1.65\,GW gross \citep{Folz2019}: 8{,}000\,EUR/kW in 2015 euros, and with the euro-area HICP (2025 over 2015: 128.75 over 100.00, factor 1.2875 \citep{EurostatHICP}) about 10{,}300 in 2025 euros; it took about 17 years \citep{Flamanville2024}, and its cost per kW is about 21\% above the central case (10{,}300 over 8{,}540, our calculation). Vogtle 3--4 cost close to 35 billion USD in total \citep{EnergytechVogtle}, of which more than 3.5--4 billion are financing charges \citep{IEEFAVogtle}, so about 31 billion without interest; for 2 $\times$ 1{,}117\,MWe net (design data of the units) that is 13{,}876\,USD/kW nominal; with uniform spending over 2013--2023 the CPI-U factor is 1.262, giving 17{,}516\,USD/kW in 2025 dollars and 15{,}656\,EUR/kW, 0.7\% above the 15{,}540 used. The extremes, 30 and 35 billion, give 15{,}150 and 17{,}680\,EUR/kW. With 15 years of construction the risk case has fixN 878.3\,EUR/kW per year at 5\% and 1{,}334.2 at 7\%; its results were not computed but read off the curve by linear interpolation between computed points. Construction takes 7--15 years, against an IAEA median of 5.7--6.3 years \citep{Ritchie2023}. Hinkley Point C, estimated at about 35 billion pounds of 2015 for 3.26\,GW, is not in the median.

\paragraph{Nuclear: other parameters.} Life is 80 years, as already licensed for 25 US units \citep{NRCslr}. Operation and fuel are 27\,EUR/MWh at 90\% load, consistent with fixed O\&M of 136--158\,USD/kW-yr \citep{Lazard2025} plus fuel, and 5\,EUR/MWh for decommissioning and waste, within the fee of 4.7--8.5\,EUR/MWh estimated in Sweden \citep{Riksgalden2025}. Stand-alone, the cost of a nuclear MWh is $\mathrm{lcoe}=K\phi\,\mathrm{CRF}(r,80)/(0.9\times8760)\times1000+32$. SMR cases use a 5-year build and 60-year life with the contract costs of Darlington (17{,}200\,EUR/kW for unit 1, 11{,}700 for four units \citep{POWERDarlington2025}) and two values that are not contracts: 6{,}250\,EUR/kW, the midpoint of the 5{,}000--7{,}500 range given for Rolls-Royce \citep{WNNWylfa2025}, and 3{,}500, a press estimate \citep{Sole24SMR}.

\paragraph{Long-duration and hydrogen storage.} The 100-hour storage costs 20\,USD/kWh over 20 years at 5\%, 1{,}486\,EUR per MWh of capacity per year with 50\% efficiency \citep{FormTarget2021,FormRTE2026}; life 20 years and O\&M of 20\% are our assumptions. It is the developer's target, not an observed price: no iron-air plant is in commercial operation and no current cost is disclosed \citep{FormMIT2026}. Hydrogen is modelled with electrolysis, geological storage and turbines. Large-scale geological hydrogen storage exists only as studies and pilots, so its costs are first-plant (PI: electrolyser 1{,}819\,EUR/kW, storage 113\,EUR/MWh-yr, turbine 1{,}079\,EUR/kW) or target (OB: 450, 56 and 1{,}324) values from studies \citep{HystoriesD86,BaringaDESNZ2025,IEAGHR2025}, not observed prices. Lithium-ion batteries are priced at the 2025 MACSE auction, an observed contract price. Imports are priced at 115.94\,EUR/MWh, the 2025 average PUN \citep{GME2025}, under the rule described in Section~\ref{sec:model}; in the variant of the earlier version the price is raised by 17\% (our assumption) in the three critical weeks of each year, with availability and credit of 38\%. The gas price in the NZ variants is 38.5\,EUR/MWh, the 2025 IG Index of the same report \citep{GME2025}.

\begin{table}[htbp]
\centering\footnotesize
\caption{Model inputs and their status. DATA: read in a source; DERIVED: calculation on data; ESTIMATE: author's judgement with a declared method; ASSUMPTION: model choice without a source, tested in a sensitivity.}\label{tab:inputs}
\input{tab_inputs}
\end{table}

\subsection{Time step, weather years and checks}\label{sec:stepmethods}
The six-hour step underestimates the value of nuclear (Section~\ref{sec:step}). Central values are therefore three-hour results. For the PI case we add a Richardson extrapolation of the gap towards a one-hour step. We assume $\mathrm{gap}(h)=G_0-c\,h^p$ with $h$ in hours; from the 12-, 6- and 3-hour gaps, $2^p$ is the ratio of the increments 12 to 6 and 6 to 3 hours, $c\,3^p=(\text{increment }6\to3)/(2^p-1)$, $G_0=\mathrm{gap}(3)+c\,3^p$ and $\mathrm{gap}(1)=G_0-c$. The extrapolation is accepted when the ratio of the increments lies between 0.35 and 0.65.
To test years outside the base sample, we rank the 38 PECD years 1982--2019 by the mean of the three worst weeks of the combined index of 171\,GW of solar and 44.6\,GW of onshore wind. The four worst years are 2015, 1998, 2014 and 1994 (13.58 to 14.66\,GW against 16.3 to 19.1\,GW for 2016--2019). They replace the weather of 2016--2019, paired in that order with the load of those years, at six hours. The criterion ignores cold spells and hydropower, and old weather is paired with recent demand structure.
\FloatBarrier

\section{Results}\label{sec:results}

\subsection{Seasonal mismatch}\label{sec:monthly}
In 2024 Italian solar plants produced 1{,}541\,GWh in January and 4{,}700\,GWh in July; wind produced 2{,}814\,GWh in January and 916\,GWh in August \citep{GSE2025,Terna2025m} (Fig.~\ref{fig:profiles}). Monthly costs in this subsection are per MWh of total demand (Section~\ref{sec:methods}). Covering the worst month with solar alone requires 1.95 times the capacity sized on the annual mean, so the real lifetime cost of 42\,EUR/MWh derived from the main FER~X auction becomes 82\,EUR/MWh of useful energy before any storage (96 at the NZIA price). Adding daily batteries at the MACSE price gives 110\,EUR/MWh (125). Covering the seasonal deficit, 16.8\% of annual energy, with batteries that cycle once a year costs about 2{,}000\,EUR/MWh; with hydrogen stored in salt caverns at the cost estimated for purely seasonal use \citep{DCU2026}, about 390\,EUR/MWh, a lower bound because the cost of electrolysers and fuel cells is excluded. A 50/50 mix of solar and wind reduces the overbuild to 1.35 and the seasonal deficit to 5\%, for 76--92\,EUR/MWh with daily batteries (84 at the reference wind price), or 90--106 with a grid cost of 14\,EUR/MWh. On this monthly view, without grid costs, the mix is cheaper than nuclear at the Dukovany price financed at 5\% (<<LC:8540:5:7|0>>\,EUR/MWh at a 90\% load factor, Table~\ref{tab:lcoe}). The monthly figures, however, ignore the hours and weeks without wind and sun \citep{KittelSchill2024} and heat waves, during which wind output in southern Europe can fall by up to 30.8\% while demand rises \citep{Molina2023}; a monitoring of large solar parks by the University of \'Evora, reported by the press, measured an average loss of 17.5\% of output over 24 hours in extreme heat, because module efficiency falls above 25\,$^\circ$C \citep{Evora2026}. The three-hour model includes them for the weather years considered, and the comparison changes.

\begin{figure}[t]
\centering\includegraphics[width=\linewidth]{fig1_profiles.pdf}
\caption{Monthly output of solar PV and wind in Italy in 2024, as a share of each technology's annual output, against flat demand. Data: \citet{GSE2025}, \citet{Terna2025m}.}\label{fig:profiles}
\end{figure}

\subsection{One curve: from fixed cost to nuclear capacity}\label{sec:least}
The central result is one curve (Fig.~\ref{fig:share}): the annualised fixed cost of nuclear, fixN in EUR/kW-yr, against the nuclear capacity of the least-cost system under net self-sufficiency (Section~\ref{sec:model}), at a 5\% and a 7\% cost of capital, for the PI and OB cost cases; the curve of the earlier import rule (imports with a capacity credit and no exports) is shown as a variant. Cost cases are points on the curve, not separate runs: each is converted to fixN with the formula of Section~\ref{sec:model} and read on the curve (Table~\ref{tab:cases}). The curve points are six-hour runs; the three-hour central points are shown as separate markers. Costs are per MWh of demand not met by existing hydropower.

\begin{figure}[t]
\centering\includegraphics[width=0.95\linewidth]{fig2_share_vs_fixedcost.pdf}
\caption{Optimal nuclear capacity against the annualised fixed cost of nuclear (fixN) under net self-sufficiency, at 5\% (blue) and 7\% (orange). Solid lines with circles: first plant (PI); dashed with squares: target (OB, reduced model). Thin grey lines: the earlier rule, imports with a capacity credit, at 5\%. Vertical lines: the five cost cases at 5\% and the central case at 7\%. Lines are six-hour runs; diamonds and crosses are the three-hour central points. Source: \texttt{results/v7.json} and \texttt{v7\_sweep.json}.}\label{fig:share}
\end{figure}

\begin{table}[htbp]
\centering
\caption{The cost cases as points on the curve, six-hour step: optimal nuclear capacity in GW and, in parentheses, share of total demand in \%, under net self-sufficiency and, as a variant, with imports with a capacity credit. fixN is 669.5 at 5\% and 922.7 at 7\% for Dukovany (runs at 670 and 923), 878.3 at 5\% and 1{,}334.2 at 7\% for the risk case. * interpolated between two computed points, or 0 beyond the last point with nuclear.}\label{tab:cases}
\input{tab_cases}
\end{table}

Under net self-sufficiency at 5\%, nuclear falls from <<N6:PI:s:5:335|1>>\,GW at Saeul cost (fixN 335, PI) to <<N6:PI:s:5:598|1>>\,GW at 598, to <<N6:PI:s:5:604|1>>\,GW at the EPC-only price (604), to <<RUN:ss0_PI_670_h6:N_GW|1>>\,GW in the central case (670) and to <<N6:PI:s:5:750|1>>\,GW at 750; a tail between 12 and 9\,GW follows up to about 880, and nuclear disappears between 950 and 1{,}050. The slope is steepest between 600 and 750 (about <<EXPR:(V('N6:PI:s:5:598')-V('N6:PI:s:5:750'))/152|2>>\,GW per EUR/kW-yr, against <<EXPR:(V('N6:PI:s:5:335')-V('N6:PI:s:5:598'))/263|2>> between 335 and 598): this is where the central case lies, and a 10\% error on construction cost moves the result by several GW. The OB curve falls earlier: <<N6:OB:s:5:750|1>>\,GW at 750 against <<N6:PI:s:5:750|1>> in PI, and zero from 950, because in OB the technologies at target cost lower the cost of the system without nuclear (<<RUN:ss0_OB_none_h6:cost_new|2>> against <<RUN:ss0_PI_none_h6:cost_new|2>>\,EUR/MWh). At 670 the two series give <<RUN:ss0_PI_670_h6:N_GW|1>> and <<RUN:ss0_OB_670_h6:N_GW|1>>\,GW. The risk case (9{,}680\,EUR/kW, 15 years, fixN 878.3) gives <<N6:PI:s:5:878|1>>\,GW in PI and <<N6:OB:s:5:878|1>> in OB; Vogtle gives none. Fifteen years of construction alone, at the central cost (fixN <<FIX:8540:5:15|1>>), give <<N6:PI:s:5:792.5|1>>\,GW in PI and <<N6:OB:s:5:792.5|1>>\,GW in OB at six hours, interpolated. Barakah (4{,}670\,EUR/kW, 8 years, fixN <<FIX:4670:5:8|1>>) gives <<N6:PI:s:5:442|1>>\,GW in PI and <<N6:OB:s:5:442|1>> in OB. Under the earlier import rule the curve is lower in the central region: <<RUN:c3_PI_670_h6:N_GW|1>> and <<RUN:c3_OB_670_h6:N_GW|1>>\,GW at 670.

\subsection{Thresholds}
Table~\ref{tab:thr} gives the fixN at which nuclear falls below 10\,GW, by linear interpolation between the two computed points that bracket 10\,GW, and the equivalent construction cost with 7 and 15 years of construction. Under net self-sufficiency at 5\% and a six-hour step the PI threshold is <<THR6:PI:s:5:10|0>> (<<EXPR:V('CPX:'+str(V('THR6:PI:s:5:10'))+':5:7')|0>>\,EUR/kW at 7 years) and the OB threshold <<THR6:OB:s:5:10|0>> (<<EXPR:V('CPX:'+str(V('THR6:OB:s:5:10'))+':5:7')|0>>). A construction cost of about 10{,}000\,EUR/kW, the range of Flamanville, keeps PI above 10\,GW and brings OB below. The central case lies <<EXPR:V('THR6:PI:s:5:10')-670|0>>\,EUR/kW-yr below the PI threshold, about <<EXPR:V('CPX:'+str(V('THR6:PI:s:5:10'))+':5:7')-8540|-2>>\,EUR/kW of construction cost, and <<EXPR:V('THR6:OB:s:5:10')-670|0>> below the OB one. At three hours the PI threshold rises: PI still has <<RUN:ss0_PI_850_h3:N_GW|1>>\,GW at fixN 850 (11{,}500\,EUR/kW at seven years), so the threshold lies beyond 850, against <<THR6:PI:s:5:10|0>> at six hours. For OB the three-hour points do not pin it down: OB has <<RUN:ss0_OB_670_h3:N_GW|1>>\,GW at 670 and <<RUN:ss0_OB_780_h3:N_GW|1>> at 780 (<<N6:OB:s:5:780|1>> at six hours), so the threshold lies between 670 and 780 as at six hours, and it cannot be said to rise. Between fixN 780 and 850 three-hour nuclear barely changes, from <<RUN:ss0_PI_780_h3:N_GW|1>> to <<RUN:ss0_PI_850_h3:N_GW|1>>\,GW in PI and from <<RUN:ss0_OB_780_h3:N_GW|1>> to <<RUN:ss0_OB_850_h3:N_GW|1>> in OB (Table~\ref{tab:step}). With 15 years of construction the thresholds fall to about <<EXPR:V('CPX:'+str(V('THR6:PI:s:5:10'))+':5:15')|-2>> and <<EXPR:V('CPX:'+str(V('THR6:OB:s:5:10'))+':5:15')|-2>>\,EUR/kW. At 7\% the PI threshold is fixN <<THR6:PI:s:7:10|0>> (<<EXPR:V('CPX:'+str(V('THR6:PI:s:7:10'))+':7:7')|0>>\,EUR/kW at 7 years): the central case, 8{,}540, lies just below it and has <<RUN:ss0_PI_7pc_923_h3:N_GW|1>>\,GW in PI and <<RUN:ss0_OB_7pc_923_h3:N_GW|1>> in OB at three hours.

\begin{table}[htbp]
\centering
\caption{fixN at which optimal nuclear capacity falls below 10\,GW, by import rule, and equivalent construction cost with 7 and 15 years of construction. * three-hour OB threshold under the capacity-credit rule: linear extrapolation, a lower bound.}\label{tab:thr}
\input{tab_thr}
\end{table}

\subsection{The time step and the central values}\label{sec:step}
Going from a six-hour to a three-hour step raises the cost of the system without nuclear (by <<EXPR:V('RUN:ss0_PI_none_h3:cost_new')-V('RUN:ss0_PI_none_h6:cost_new')|1>>\,EUR/MWh in PI) more than that of the system with nuclear (<<EXPR:V('RUN:ss0_PI_670_h3:cost_new')-V('RUN:ss0_PI_670_h6:cost_new')|1>>), so the gap grows and, at equal fixN, so does optimal nuclear capacity (Table~\ref{tab:step}). A coarse step averages out short-term variability that solar, wind and batteries must cover and nuclear baseload need not follow; this explanation is plausible but not verified. The six-hour step therefore underestimates the value of nuclear and its optimal capacity at every point computed at three hours (the threshold rises for certain only in PI), and absolute six-hour cost levels are underestimated by an amount that cannot be extrapolated: at two years (run only at fixN 598) the cost of PI rises by 1.60\,EUR/MWh from 12 to 6 hours and by 2.49 from 6 to 3 hours. The gap converges better, and is the decision metric.

\begin{table}[htbp]
\centering
\caption{Costs and capacities at six- and three-hour steps, four years, net self-sufficiency, 5\% unless stated; cost in EUR/MWh.}\label{tab:step}
\input{tab_step}
\end{table}

Central values are the three-hour results (Table~\ref{tab:central}): PI <<RUN:ss0_PI_670_h3:N_GW|1>>\,GW (<<RUN:ss0_PI_670_h3:sh|1>>\% of demand) and OB <<RUN:ss0_OB_670_h3:N_GW|1>>\,GW (<<RUN:ss0_OB_670_h3:sh|1>>\%), with gaps of <<GAP:ss0_PI_none_h3:ss0_PI_670_h3|2>> and <<GAP:ss0_OB_none_h3:ss0_OB_670_h3|2>>\,EUR/MWh (<<GAP:ss0_PI_none_h6:ss0_PI_670_h6|2>> and <<GAP:ss0_OB_none_h6:ss0_OB_670_h6|2>> at six hours); about <<EXPR:V('GAP:ss0_PI_none_h3:ss0_PI_670_h3')*356.6/1000|1>> billion euros a year for PI (<<GAP:ss0_PI_none_h3:ss0_PI_670_h3|2>> times 356.6\,TWh, our calculation). The Richardson extrapolation (Section~\ref{sec:stepmethods}) at fixN 670 with four years has an increment ratio of <<RICHS:PI:670:ratio|2>> for PI, within the acceptable range of 0.35--0.65, and gives a one-hour gap of <<RICHS:PI:670:G1|2>>\,EUR/MWh against <<RICHS:PI:670:g3|2>> at three hours; for OB the ratio is <<RICHS:OB:670:ratio|2>>; under the earlier import rule the PI ratio was <<RICH:PI:670:ratio|2>>, with a one-hour gap of <<RICH:PI:670:G1|2>>\,EUR/MWh against <<RICH:PI:670:g3|2>> at three hours. Where the ratio is outside the range the extrapolation is not reliable: the sign of the correction is certain, its size is not, and the three-hour gap is a lower bound (Table~\ref{tab:conv}). The headline figures use the three-hour values.

\begin{table}[htbp]
\centering
\caption{Central values under net self-sufficiency, and the variant with imports with a capacity credit (earlier rule). Cost in EUR/MWh; share of total demand in \%; gap = cost without nuclear minus cost with nuclear. The central value is the three-hour result.}\label{tab:central}
\input{tab_central}
\end{table}

\begin{table}[htbp]
\centering
\caption{Convergence of the gap with the time step (EUR/MWh). The limit for $h\to0$ is $G_0$ of the Richardson estimate. * ratio outside 0.35--0.65: extrapolation not reliable (n.a., not computed). Rows under the capacity-credit rule at fixN 598 (two and four years) and 670 (four years), and under net self-sufficiency at 670 (four years).}\label{tab:conv}
\input{tab_conv}
\end{table}

\subsection{The price of import dependence, and 7\%}\label{sec:dep}
Table~\ref{tab:dep} compares import rules at six hours, with nuclear at the central cost and without nuclear. The relevant comparison is not between today's import share and that of 2040, but between a system that does not rely on others and one that accepts a fixed share of their energy. With nuclear, allowing net imports of up to 10\% of demand lowers the system cost from <<RUN:ss0_PI_670_h6:cost_new|2>> to <<RUN:ss10_PI_670_h6:cost_new|2>>\,EUR/MWh in PI and from <<RUN:ss0_OB_670_h6:cost_new|2>> to <<RUN:ss10_OB_670_h6:cost_new|2>> in OB, by <<EXPR:V('RUN:ss0_PI_670_h6:cost_new')-V('RUN:ss10_PI_670_h6:cost_new')|2>> and <<EXPR:V('RUN:ss0_OB_670_h6:cost_new')-V('RUN:ss10_OB_670_h6:cost_new')|2>>\,EUR/MWh, and the constraint does not even bind: the system imports <<EXPR:V('UN:ss10_PI_670_h6:IMPORT:TWh')-V('UN:ss10_PI_670_h6:EXPORT:TWh')|0>>\,TWh net in PI out of 40 allowed. Without nuclear the same allowance is worth <<EXPR:V('RUN:ss0_PI_none_h6:cost_new')-V('RUN:ss10_PI_none_h6:cost_new')|2>> (PI) and <<EXPR:V('RUN:ss0_OB_none_h6:cost_new')-V('RUN:ss10_OB_none_h6:cost_new')|2>>\,EUR/MWh (OB): it is the system without nuclear that needs other countries' energy. A 10\% net import target thus buys about 1--1.5\% of system cost with a structural dependence. The earlier rule is cheaper still, <<RUN:c3_PI_670_h6:cost_new|2>>\,EUR/MWh in PI, because it relies on about 7\,GW of foreign capacity in the critical weeks (38\% of the import capacity) and never returns the energy it takes; that value comes from capacity that other countries build and pay for, and would not exist if every country planned the same way.

\begin{table}[htbp]
\centering
\caption{The price of import dependence: system cost at six hours, four years, PI and OB, with nuclear at fixN 670 and without nuclear, by import rule. $X$ is the allowed annual net import as a share of demand. Net imports are imports minus exports in the case with nuclear. Cost in EUR/MWh.}\label{tab:dep}
\input{tab_dep}
\end{table}

At a 7\% cost of capital the central case (fixN 923) gives <<RUN:ss0_PI_7pc_923_h6:N_GW|1>>\,GW in PI and <<RUN:ss0_OB_7pc_923_h6:N_GW|1>> in OB at six hours, <<RUN:ss0_PI_7pc_923_h3:N_GW|1>> and <<RUN:ss0_OB_7pc_923_h3:N_GW|1>> at three hours, with gaps of <<GAP:ss0_PI_7pc_none_h3:ss0_PI_7pc_923_h3|2>> and <<GAP:ss0_OB_7pc_none_h3:ss0_OB_7pc_923_h3|2>>\,EUR/MWh at three hours. At 7\% the costs of solar, wind, storage and hydrogen also rise, so at equal fixN the 7\% curve lies above the 5\% one; it is the fixN of nuclear that rises most, from 670 to 923. The risk case at 7\% (fixN 1{,}334.2) gives no nuclear; nuclear is zero from 1{,}100. Without imports of any kind, a stricter variant of self-sufficiency, the central case gives <<N6:PI:n:5:670|1>>\,GW in PI and <<N6:OB:n:5:670|1>> in OB at 5\%, and about <<N6:PI:n:7:923|1>> and <<N6:OB:n:7:923|1>> at 7\% (six hours, interpolated). The full points are in Tables~\ref{tab:curve5} and~\ref{tab:other}.

\subsection{The central system}
At three hours the central PI system (fixN 670) has <<UN:ss0_PI_670_h3:NUC:GW|1>>\,GW of nuclear (<<UN:ss0_PI_670_h3:NUC:hours|0>> equivalent hours, <<UN:ss0_PI_670_h3:NUC:TWh|1>>\,TWh), <<GR:ss0_PI_670_h3:pv:GW|1>>\,GW of solar, <<GR:ss0_PI_670_h3:on:GW|1>>\,GW of onshore wind, <<UN:ss0_PI_670_h3:IMPORT:GW|1>>\,GW of interconnection (<<UN:ss0_PI_670_h3:IMPORT:TWh|1>>\,TWh imported and as much exported) and <<GR:ss0_PI_670_h3:battery:GWh|0>>\,GWh of batteries. At six hours it has <<UN:ss0_PI_670_h6:NUC:GW|1>>\,GW of nuclear (<<UN:ss0_PI_670_h6:NUC:hours|0>>\,h, <<UN:ss0_PI_670_h6:NUC:TWh|1>>\,TWh), <<GR:ss0_PI_670_h6:pv:GW|1>>\,GW of solar (24.2 existing, <<EXPR:V('GR:ss0_PI_670_h6:pv:GW')-V('UN:ss0_PI_670_h6:PVex_terra:GW')-V('UN:ss0_PI_670_h6:PVex_tetti:GW')|1>> new: <<EXPR:V('UN:ss0_PI_670_h6:PV_terra_Sud-isole:GW')+V('UN:ss0_PI_670_h6:PV_terra-pipeline_Lazio/Sicilia/Piem./Puglia/Sard.:GW')+V('UN:ss0_PI_670_h6:PV_terra_Centro-Sud:GW')+V('UN:ss0_PI_670_h6:PV_terra_Nord:GW')|1>> ground-mounted, all ground blocks exhausted and <<UN:ss0_PI_670_h6:PV_agrivoltaico_Sud-isole:GW|1>> agrivoltaic), <<GR:ss0_PI_670_h6:on:GW|1>>\,GW of onshore wind (E0 13.6, E1 6.0, E2 <<UN:ss0_PI_670_h6:E2:GW|1>>), <<UN:ss0_PI_670_h6:IMPORT:GW|1>>\,GW of interconnection (<<UN:ss0_PI_670_h6:IMPORT:TWh|1>>\,TWh imported and as much exported) and <<GR:ss0_PI_670_h6:battery:GWh|0>>\,GWh of batteries; solar and wind curtailment is about <<EXPR:(FULL['ss0_PI_670_h6']['diag']['VRE']['on']['curt']+FULL['ss0_PI_670_h6']['diag']['VRE']['pv']['curt'])/(FULL['ss0_PI_670_h6']['diag']['VRE']['on']['avail']+FULL['ss0_PI_670_h6']['diag']['VRE']['pv']['avail'])*100|1>>\% of available output. Its cost of <<RUN:ss0_PI_670_h6:cost_new|2>>\,EUR/MWh comprises nuclear <<COMP:ss0_PI_670_h6:nuclear|2>>, solar <<COMP:ss0_PI_670_h6:pv|2>>, wind <<COMP:ss0_PI_670_h6:on|2>>, imports <<COMP:ss0_PI_670_h6:import|2>>, batteries <<COMP:ss0_PI_670_h6:battery|2>>, geothermal <<COMP:ss0_PI_670_h6:geo|2>>, bioenergy <<COMP:ss0_PI_670_h6:bio|2>>, 100-hour storage <<COMP:ss0_PI_670_h6:ldes|2>> and interruptible load <<COMP:ss0_PI_670_h6:interr|2>>. The central OB system at six hours has <<UN:ss0_OB_670_h6:NUC:GW|1>>\,GW of nuclear, <<GR:ss0_OB_670_h6:pv:GW|1>>\,GW of solar, <<GR:ss0_OB_670_h6:on:GW|1>>\,GW of onshore wind, <<GR:ss0_OB_670_h6:battery:GWh|0>>\,GWh of batteries and <<GR:ss0_OB_670_h6:caes:GWh|0>>\,GWh of compressed air (Fig.~\ref{fig:capacity}).

\begin{figure}[t]
\centering\includegraphics[width=0.95\linewidth]{fig3_capacity.pdf}
\caption{Installed capacity and cost of the least-cost system at a six-hour step for PI and OB, with nuclear at fixN 670 (5\%), at 923 (7\%) and without nuclear, under net self-sufficiency. PV includes the existing fleet after 2016; onshore wind includes E0.}\label{fig:capacity}
\end{figure}

\begin{figure}[t]
\centering\includegraphics[width=0.7\linewidth]{fig4_cost_vs_nuclear.pdf}
\caption{Gap (cost without nuclear minus cost with nuclear) against fixN under net self-sufficiency, at 5\% and 7\%. Lines are six-hour runs; isolated points are the three-hour central values (PI <<GAP:ss0_PI_none_h3:ss0_PI_670_h3|2>> and OB <<GAP:ss0_OB_none_h3:ss0_OB_670_h3|2>>\,EUR/MWh at 670; PI <<GAP:ss0_PI_7pc_none_h3:ss0_PI_7pc_923_h3|2>> and OB <<GAP:ss0_OB_7pc_none_h3:ss0_OB_7pc_923_h3|2>> at 7\%, fixN 923).}\label{fig:gap}
\end{figure}

\subsection{The system without nuclear and the bridge from the previous version}\label{sec:bridge}
Under net self-sufficiency, at six hours the PI system without nuclear has <<GR:ss0_PI_none_h6:pv:GW|0>>\,GW of solar, <<GR:ss0_PI_none_h6:on:GW|1>>\,GW of onshore wind, <<UN:ss0_PI_none_h6:IMPORT:GW|1>>\,GW of interconnection (<<UN:ss0_PI_none_h6:IMPORT:TWh|1>>\,TWh imported and as much exported), <<GR:ss0_PI_none_h6:battery:GWh|0>>\,GWh of batteries, <<GR:ss0_PI_none_h6:h2:GWh|0>>\,GWh of hydrogen storage with a <<UN:ss0_PI_none_h6:H2:Pd_GW|1>>\,GW turbine and <<GR:ss0_PI_none_h6:ldes:GWh|0>>\,GWh of 100-hour storage; OB has <<GR:ss0_OB_none_h6:pv:GW|0>>\,GW of solar, <<GR:ss0_OB_none_h6:on:GW|1>>\,GW of onshore wind, <<GR:ss0_OB_none_h6:caes:GWh|0>>\,GWh of compressed air, <<GR:ss0_OB_none_h6:h2:GWh|0>>\,GWh of hydrogen storage and <<GR:ss0_OB_none_h6:battery:GWh|0>>\,GWh of batteries. It costs <<RUN:ss0_PI_none_h6:cost_new|2>> (PI) and <<RUN:ss0_OB_none_h6:cost_new|2>>\,EUR/MWh (OB) at six hours, <<RUN:ss0_PI_none_h3:cost_new|2>> and <<RUN:ss0_OB_none_h3:cost_new|2>> at three. The previous version of this paper, whose model lacked these technologies, gave 204.2\,EUR/MWh and 38.5\,GW of nuclear supplying 59\% of demand. Table~\ref{tab:bridgepi} rebuilds the difference one technology at a time at six hours and with nuclear at fixN 598, the cost of the previous version, which is not the present central case (670); step 0 reproduces the previous version in the present code (193.01\,EUR/MWh without nuclear and 94.05 with nuclear at fixN 598, 38.4\,GW; 204.2 is the same model at three hours).

\begin{table}[htbp]
\centering
\caption{From the previous model to the present one, PI case, six hours, cumulative; cost in EUR/MWh. The last column is the change in cost of the system with nuclear. The OB case is Table~\ref{tab:bridgeob}.}\label{tab:bridgepi}
\input{tab_bridge_pi}
\end{table}

Imports explain most of the reduction without nuclear: 72.67 of 103.20\,EUR/MWh (193.01 minus 89.81). Hydrogen with electrolysis, geological storage and turbine removes 16.00, bioenergy 7.81, geothermal 5.11, smart charging and interruptible load 0.59 and 100-hour storage 1.00; concentrated solar power, compressed air and new pumped hydro are worth nothing in PI, while in OB the first is worth 1.91 and the second 3.80. With nuclear at fixN 598 (as in the table), nuclear falls from 38.4\,GW to 27.9 with imports, 26.4 with geothermal and 23.3 with bioenergy, and rises to 23.9 with smart charging; hydrogen does not change optimal nuclear. The bridge uses the earlier import rule (capacity credit, no exports); the last step, from that rule to net self-sufficiency, raises the cost without nuclear from 89.81 to <<RUN:ss0_PI_none_h6:cost_new|2>>\,EUR/MWh and that with nuclear at fixN 598 from <<RUN:sw_PI_imp_598:cost_new|2>> to <<RUN:ss0_PI_598_h6:cost_new|2>>, and nuclear from <<RUN:sw_PI_imp_598:N_GW|1>> to <<RUN:ss0_PI_598_h6:N_GW|1>>\,GW (PI, six hours). Large-scale geological hydrogen storage exists only as studies and pilots, and the hydrogen-to-power chain has no commercial plant: model costs are first-plant or target values from studies, not observed prices. Without nuclear hydrogen is worth 16\,EUR/MWh, the most uncertain part of the cost of the system without nuclear; the gap is computed with those costs.

\subsection{Sensitivity analyses}\label{sec:sens}
The sensitivities of this subsection were computed under the earlier import rule (capacity credit, no exports) and were not repeated under net self-sufficiency. They are compared with the central case of the same rule at six hours: <<RUN:c3_PI_670_h6:N_GW|1>>\,GW and a gap of <<GAP:val_PI_none:c3_PI_670_h6|2>>\,EUR/MWh in PI, <<RUN:c3_OB_670_h6:N_GW|1>>\,GW and <<GAP:val_OB_none:c3_OB_670_h6|2>> in OB. Under net self-sufficiency the base gap is higher (<<GAP:ss0_PI_none_h6:ss0_PI_670_h6|2>>\,EUR/MWh in PI and <<GAP:ss0_OB_none_h6:ss0_OB_670_h6|2>> in OB), but the effect of each variation under the new rule was not recomputed and its size is unknown: CSP, gas with DAC or SMRs could matter more, because they would also replace the balancing that imports provided before.

\paragraph{Concentrated solar power.} CSP enters only in OB and, in the base case, only without nuclear; its 10\,GW ceiling is our assumption. With nuclear at fixN 670 the OB gap is <<GAP:val_OB_none:c3_OB_670_h6|2>>\,EUR/MWh; a 3\,GW ceiling gives <<GAP:se_OB_csp_c3_none:c3s_se_OB_csp_c3_598|2>>, a yield of 2{,}800 h gives <<GAP:se_OB_csp_y2800_none:c3s_se_OB_csp_y2800_598|2>>, 3{,}300 h gives <<GAP:se_OB_csp_y3300_none:c3s_se_OB_csp_y3300_598|2>> and an 8-hour capacity credit <<GAP:se_OB_csp_ch8_none:c3s_se_OB_csp_ch8_598|2>> (Table~\ref{tab:sens}). With nuclear, CSP enters only with the 3{,}300 h yield, and then it lowers optimal nuclear from <<RUN:c3_OB_670_h6:N_GW|1>> to <<RUN:c3s_se_OB_csp_y3300_598:N_GW|1>>\,GW (six hours): at fixN 670 nuclear is more sensitive to the CSP yield than at 598. In the other three cases nuclear changes by at most <<EXPR:max(abs(V('RUN:c3s_se_OB_csp_'+k+'_598:N_GW')-V('RUN:c3_OB_670_h6:N_GW')) for k in ('c3','y2800','ch8'))|2>>\,GW.

\paragraph{Imports at double price.} Doubling the price of imports in the critical weeks, instead of the 17\% mark-up, gives gaps of <<GAP:se_PI_pm2_none:c3s_se_PI_pm2_598|2>> (PI) and <<GAP:se_OB_pm2_none:c3s_se_OB_pm2_598|2>>\,EUR/MWh (OB), against <<GAP:val_PI_none:c3_PI_670_h6|2>> and <<GAP:val_OB_none:c3_OB_670_h6|2>>, and optimal nuclear of <<RUN:c3s_se_PI_pm2_598:N_GW|1>>\,GW (PI) and <<RUN:c3s_se_OB_pm2_598:N_GW|1>> (OB): the value of nuclear is insensitive to that price.

\begin{table}[htbp]
\centering
\caption{Sensitivities, six hours, four years, imports with a capacity credit (earlier rule); cost in EUR/MWh.}\label{tab:sens}
\input{tab_sens}
\end{table}

\paragraph{Small modular reactors.} With large nuclear switched off, SMRs are not built at contract costs: Darlington's 17{,}200 and 11{,}700\,EUR/kW correspond to fixN 1{,}179 and 850, above the thresholds. They enter only at two values that are not contracts, 6{,}250 and 3{,}500\,EUR/kW, with 27 and 38\,GW and system costs of 74.78 and 59.20\,EUR/MWh in PI (Table~\ref{tab:smr}). These are developers' and press figures, not offers: boundary scenarios, not results.

\begin{table}[htbp]
\centering
\caption{Small modular reactors with large nuclear off, six hours, four years, imports with a capacity credit. SMR fixN with a 5-year build, 60-year life, $r=5\%$.}\label{tab:smr}
\input{tab_smr}
\end{table}

\paragraph{Net-zero emissions.} In the NZ mode gas with CCS and direct air capture (DAC) are allowed. With nuclear at fixN 670 gas is not chosen in the cases tried (first-plant costs). Without nuclear, gas with CCS does not enter at first-plant costs; at target costs it enters in PI with 6.5\,GW (cost 88.68). Gas without capture with all emissions offset by DAC lowers the cost without nuclear to 85.00 in PI (13.4\,GW of gas, 27.6\,Mt CO$_2$ a year of DAC at target cost) and 83.14 in OB (10.3\,GW, 19.1\,Mt). With target DAC the gap falls from <<GAP:val_PI_none:c3_PI_670_h6|2>> to <<GAP:se_NZ_PI_none_noccs_dactarget:c3s_se_NZ_PI_nuc598|2>> in PI and from <<GAP:val_OB_none:c3_OB_670_h6|2>> to <<GAP:se_NZ_OB_none_noccs:c3s_se_NZ_OB_nuc598|2>> in OB (six hours, against ZD nuclear at fixN 670). The result depends on a DAC cost of 100\,EUR/t (target) or 650 (first plant), our assumptions; 27.6\,Mt a year is about 770 times the nameplate of the largest plant in the world, 36{,}000\,t a year \citep{Climeworks2026}, so the case marks the limit of the assumption favourable to gas rather than an estimate.

\begin{table}[htbp]
\centering
\caption{NZ variants, six hours, four years, imports with a capacity credit.}\label{tab:nz}
\input{tab_nz}
\end{table}

\subsection{Worst weather years}
The four worst PECD years are 2015, 1998, 2014 and 1994 (mean of the three worst weeks of the solar and wind index 13.58 to 14.66\,GW against 16.3 to 19.1 for 2016--2019). Under net self-sufficiency, with them the cost without nuclear rises by <<EXPR:V('RUN:ss0_wx4_PI_none_worst:cost_new')-V('RUN:ss0_PI_none_h6:cost_new')|2>>\,EUR/MWh (PI) and <<EXPR:V('RUN:ss0_wx4_OB_none_worst:cost_new')-V('RUN:ss0_OB_none_h6:cost_new')|2>> (OB), that with nuclear by <<EXPR:V('RUN:ss0_wx4_PI_670_worst:cost_new')-V('RUN:ss0_PI_670_h6:cost_new')|2>> and <<EXPR:V('RUN:ss0_wx4_OB_670_worst:cost_new')-V('RUN:ss0_OB_670_h6:cost_new')|2>>, and optimal nuclear by <<EXPR:V('RUN:ss0_wx4_PI_670_worst:N_GW')-V('RUN:ss0_PI_670_h6:N_GW')|1>> and <<EXPR:V('RUN:ss0_wx4_OB_670_worst:N_GW')-V('RUN:ss0_OB_670_h6:N_GW')|1>>\,GW, to <<RUN:ss0_wx4_PI_670_worst:N_GW|0>>\,GW (<<RUN:ss0_wx4_PI_670_worst:sh|0>>\%) in PI and <<RUN:ss0_wx4_OB_670_worst:N_GW|0>> (<<RUN:ss0_wx4_OB_670_worst:sh|0>>\%) in OB. The gap grows to <<EXPR:V('RUN:ss0_wx4_PI_none_worst:cost_new')-V('RUN:ss0_wx4_PI_670_worst:cost_new')|2>> against <<GAP:ss0_PI_none_h6:ss0_PI_670_h6|2>> in PI and <<EXPR:V('RUN:ss0_wx4_OB_none_worst:cost_new')-V('RUN:ss0_wx4_OB_670_worst:cost_new')|2>> against <<GAP:ss0_OB_none_h6:ss0_OB_670_h6|2>> in OB (Table~\ref{tab:wx}). Nuclear protects against low-wind low-sun periods, which are less severe in the four base years; without help from abroad the system without nuclear covers them with storage and capacity used a few weeks a year.

\begin{table}[htbp]
\centering
\caption{Worst weather years against the base years, six hours, net self-sufficiency, nuclear at fixN 670; cost in EUR/MWh.}\label{tab:wx}
\input{tab_wx}
\end{table}
\FloatBarrier

\section{Discussion}\label{sec:discussion}

\subsection{Why other studies reach different conclusions}
The results differ from widely cited cost comparisons for identifiable reasons. \citet{Lazard2025} reports 175--255\,USD/MWh for US nuclear, derived from Vogtle 3 and 4, the first US units in thirty years, at 12{,}300--17{,}570\,USD/kW and 7.7\% cost of capital: a first-of-a-kind American cost. \citet{Fraunhofer2024} reports 136--490\,EUR/MWh for Germany assuming 4{,}300--6{,}300 operating hours today and 2{,}000--4{,}000 by 2045, i.e.\ a reactor idle for more than half of the year, together with 6{,}000--16{,}000\,EUR/kW and 7.8\% real cost of capital; in our central PI system nuclear runs about <<UN:ss0_PI_670_h3:NUC:hours|0>> equivalent hours (<<EXPR:V('UN:ss0_PI_670_h3:NUC:hours')/87.6|0>>\%), and the stand-alone cost at 90\% (Table~\ref{tab:lcoe}) understates its cost per MWh slightly. The same report gives 6--22.5\,ct/kWh for solar with batteries, for systems that shift part of their own output by a few hours without any guarantee of covering load. \citet{Lund2026} find renewables 53\% cheaper than nuclear in Denmark with an effective nuclear cost of 10{,}000\,EUR/kW and 8\% cost of capital. The 188--991\,USD/MWh reported by Banca d'Italia \citep{Bankitalia2025} comes from a study of small modular reactors \citep{Steigerwald2023}, not of the large reactors considered here, and the Italia Solare scenarios assume 11{,}000\,EUR/kW in 2040 \citep{ItaliaSolare2026}, above the Dukovany contract. Their Vogtle-based figures are our risk cases, not the central case.

System studies without nuclear rely on technologies without market prices for the winter. \citet{Lombardi2020} cover low-wind, low-sun periods with existing gas plants burning synthetic methane from electrolysis and direct air capture, which would require about 3.5 billion m$^3$ per year (our estimate) against 0.4--0.57 billion m$^3$ of biomethane produced today \citep{Biometano2024}. \citet{Franzmann2026} obtain 91.6\,EUR/MWh for a fully renewable Europe with 1{,}116\,GW of electrolysers and 501\,GW of hydrogen turbines, hydrogen accounting for 22\% of system cost, without sensitivity analysis on hydrogen or grid costs. The study closest to ours, \citet{Mastrantuono2026}, finds nuclear competitive, but its 8\,GWe cap binds; our model, without a cap, at that cost (about 3{,}600\,EUR/kW at 1.1188\,USD per euro, without inflation, with 8 years of construction: fixN 373 at 5\% and 488 at 7\%, our calculation) chooses about <<N6:PI:s:5:373|0>>\,GW at 5\% (PI, six hours, interpolated) and at 7\% more than the <<RUN:ss0_PI_7pc_700_h6:N_GW|0>>\,GW of the first computed point (fixN 700), far above 8\,GWe, so the cap would also bind in our model. The unpublished Politecnico di Milano analysis \citep{PoliMi2026} finds up to 15\,GW of nuclear at 5{,}500\,EUR/kW and none at 8{,}300\,EUR/kW with ten-year construction. We do not know its discount rate. With seven years, 5{,}500\,EUR/kW is fixN 485 at 5\% and 649 at 7\%, and our model gives <<N6:PI:s:5:485|0>>\,GW (PI, six hours) at 5\% and more than <<RUN:ss0_PI_7pc_700_h6:N_GW|0>> at 7\%, above their 15\,GW; with ten years, 8{,}300\,EUR/kW is fixN 696 at 5\% and 989 at 7\%, and our model gives <<N6:PI:s:5:696|1>>\,GW (PI) and <<N6:OB:s:5:696|1>> (OB) at 5\% and <<N6:PI:s:7:989|1>> and <<N6:OB:s:7:989|1>> at 7\%, so under net self-sufficiency their conclusion of no nuclear does not coincide with ours in any case. In the central PI case at three hours a system without nuclear costs <<EXPR:(V('RUN:ss0_PI_none_h3:cost_new')/V('RUN:ss0_PI_670_h3:cost_new')-1)*100|0>>\% more than the system with nuclear (<<RUN:ss0_PI_none_h3:cost_new|2>> against <<RUN:ss0_PI_670_h3:cost_new|2>>\,EUR/MWh) and in OB <<EXPR:(V('RUN:ss0_OB_none_h3:cost_new')/V('RUN:ss0_OB_670_h3:cost_new')-1)*100|0>>\% more (<<RUN:ss0_OB_none_h3:cost_new|2>> against <<RUN:ss0_OB_670_h3:cost_new|2>>): both are below the 11--163\% increase implied by \citet{Sepulveda2018} (our conversion of their 10--62\% reduction), and the gaps in EUR/MWh (<<GAP:ss0_PI_none_h3:ss0_PI_670_h3|2>> and <<GAP:ss0_OB_none_h3:ss0_OB_670_h3|2>>) are in the low range of the 8--50\,USD/MWh integration costs of \citet{NEA2019}, because the model has imports, hydrogen and other firm sources. Without imports at five per cent and six hours, with nuclear at fixN 670, the increase is <<EXPR:(120.17/V('C6:PI:n:5:670')-1)*100|0>>\% (PI, 120.17 against <<C6:PI:n:5:670|2>>) and <<EXPR:(101.57/V('C6:OB:n:5:670')-1)*100|0>>\% (OB, 101.57 against <<C6:OB:n:5:670|2>>), with costs interpolated between computed points. The comparison is approximate because systems and costs differ.

\citet{Pastore2025} describe a fully renewable Italy in 2050 with an hourly model, assuming a potential of 115.4\,GW of onshore and 55.7\,GW of offshore wind, and install the whole potential in most scenarios. Our central onshore limit is 69.6\,GW, in blocks of 1{,}200--2{,}400 full-load hours at rising cost; the JRC estimate under current siting rules is 96.9\,GW \citep{JRCEUTIMES}. Their result therefore rests on sites beyond our curve. In our model the central PI system builds <<GR:ss0_PI_670_h6:on:GW|1>>\,GW of onshore wind at six hours (<<GR:ss0_PI_670_h3:on:GW|1>> at three) and leaves about 40\,GW of blocks E3-E5 unused, at costs rising to 118\,EUR/MWh: a longer curve at rising cost would lower mainly the cost of the system without nuclear.

Models that impose net-zero rather than zero direct emissions can cover winter with gas plants fitted with carbon capture and storage and offset the residual with direct air capture. No gas-fired power plant with capture is yet in operation: the first, Net Zero Teesside in the UK (742\,MW, up to about 2\,Mt CO$_2$ a year), is expected in 2028, and its contract terms are not public \citep{NZTPower}. The only power plant with capture and years of operation, the coal-fired Boundary Dam 3 in Canada, captured on average 57\% of its CO$_2$ against a 90\% target, with capture running in 80\% of operating hours \citep{IEEFABD3}. In Italy, Ravenna CCS targets 4\,Mt a year by 2030 from industrial emitters \citep{EniRavenna2026}. The most favourable full-chain cost estimate, about GBP\,110 per tonne, adds about GBP\,37 per MWh, and other estimates are twice as high \citep{Goodall2025}. At early-October 2026 gas prices (PSV about 74\,EUR/MWh, \citealp{SelectraPSV2026}) the fuel alone of a combined cycle with capture costs about 150\,EUR per MWh of electricity (our calculation at 50\% efficiency), against <<RUN:ss0_PI_670_h3:cost_new|1>>\,EUR/MWh for the whole system with nuclear in the central case at three hours. In the model gas exists only in the NZ variants and costs 38.5\,EUR/MWh, the 2025 average PSV \citep{GME2025}, about half the autumn 2026 price, so the variant favours gas. Direct air capture is further away: the largest plant, Mammoth in Iceland, captured 675\,t in the first half of 2026 against a nameplate of 36{,}000\,t a year \citep{Climeworks2026}. A low-nuclear result that depends on these technologies holds only if they reach the assumed cost and performance. This can be checked after 2028: if Teesside shows at least 90\% capture at costs near GBP\,100 per tonne, the gap of nuclear should be recomputed with the NZ variants of Section~\ref{sec:sens}.

\subsection{Construction cost and time}
The nuclear share depends on the annualised fixed cost, hence on three variables that policy can influence. Construction cost falls with series production of one design: domestic Korean units cost 2{,}400--3{,}300\,USD/kW, the exported Barakah units 3{,}600--4{,}400\,USD/kW, the Dukovany contract, 8{,}540\,EUR/kW in the model (EPC price net of fuel, plus owner's costs) \citep{HSF2026,KHNP2025,WNAuae,Dukovany2024}. Cost of capital falls with public financing: the European Commission approved for Dukovany a subsidised state loan covering most construction costs and a 40-year contract for difference \citep{WNN2024}. Construction time falls with completed engineering and stable rules. Flamanville, at about 10{,}300\,EUR/kW in 2025 euros \citep{EurostatHICP} (8{,}000 in 2015 euros), took about 17 years \citep{Flamanville2024} and Hinkley Point C will take about 13 \citep{Hinkley2026}, but the median for reactors started since 1990 is 5.7 years and 89\% were completed in less than 10 \citep{Ritchie2023}. For Flamanville, the Folz report points to industrial causes, incomplete engineering and loss of skills \citep{Folz2019}; in the United States, rules changed during construction after Three Mile Island raised costs and durations \citep{Lovering2016}.


\subsection{Deployment rate and siting}
The central optimum, <<RUN:ss0_PI_670_h3:N_GW|1>>\,GW at three hours (PI), corresponds to about <<EXPR:V('RUN:ss0_PI_670_h3:N_GW')/1.055|0>> APR1000 units of 1{,}055\,MWe, the unit of the central case, or about <<EXPR:V('RUN:ss0_PI_670_h3:N_GW')/1.4|0>> APR1400 of 1.4\,GW. This number is a constraint to verify, not a fact: the model describes a regime system with 2040 demand and all fleets already built, not the situation in 2040. With a first unit in 2035, <<EXPR:V('RUN:ss0_PI_670_h3:N_GW')/1.055|0>> APR1000 units by 2050 require about <<EXPR:V('RUN:ss0_PI_670_h3:N_GW')/1.055/15|1>> a year for 15 years, about <<EXPR:V('RUN:ss0_PI_670_h3:N_GW')/15|1>>\,GW a year (<<RUN:ss0_PI_670_h3:N_GW|1>> over 15). According to the IAEA reactor tables, France connected 52 reactors to the grid between 1977 and 1990, 54.4\,GW net: 3.9\,GW a year on average, with peaks of 8 reactors in 1981 and 7.4\,GW in 1986 \citep{IAEAFrance}. Italy would need about <<EXPR:V('RUN:ss0_PI_670_h3:N_GW')/15/3.9*100|0>>\% of the average French pace of that period (<<EXPR:V('RUN:ss0_PI_670_h3:N_GW')/15|1>> over 3.9). Korea commissioned about one reactor a year in the 1980s \citep{WNAcountries}. The required pace is below the French and similar to the Korean one, in a country restarting without a full supply chain. The model does not compute the cost of a slower pace: if the pace stays below <<EXPR:V('RUN:ss0_PI_670_h3:N_GW')/1.055/15|1>> reactors a year the steady-state fleet is reached later, and wind, solar, imports and batteries can be added meanwhile.

Siting differs sharply between the options. Reactors are usually built in groups of four to six per site, so <<EXPR:V('RUN:ss0_PI_670_h3:N_GW')/1.055|0>> APR1000 units require <<EXPR:V('RUN:ss0_PI_670_h3:N_GW')/1.055/6|0>> to <<EXPR:V('RUN:ss0_PI_670_h3:N_GW')/1.055/4|0>> sites, which could include former thermal plants already connected to the extra-high-voltage grid if they meet seismic and cooling requirements. The central system (PI, six hours) needs <<EXPR:V('GR:ss0_PI_670_h6:on:GW')-13.6|1>>\,GW of new onshore wind, about <<EXPR:(V('GR:ss0_PI_670_h6:on:GW')-13.6)*200|-1>> turbines of 5\,MW in <<EXPR:(V('GR:ss0_PI_670_h6:on:GW')-13.6)*10|-1>>--<<EXPR:(V('GR:ss0_PI_670_h6:on:GW')-13.6)*20|-1>> parks, and <<EXPR:V('UN:ss0_PI_670_h6:PV_terra_Sud-isole:GW')+V('UN:ss0_PI_670_h6:PV_terra-pipeline_Lazio/Sicilia/Piem./Puglia/Sard.:GW')+V('UN:ss0_PI_670_h6:PV_terra_Centro-Sud:GW')+V('UN:ss0_PI_670_h6:PV_terra_Nord:GW')|1>>\,GW of new ground-mounted solar, about <<EXPR:(V('UN:ss0_PI_670_h6:PV_terra_Sud-isole:GW')+V('UN:ss0_PI_670_h6:PV_terra-pipeline_Lazio/Sicilia/Piem./Puglia/Sard.:GW')+V('UN:ss0_PI_670_h6:PV_terra_Centro-Sud:GW')+V('UN:ss0_PI_670_h6:PV_terra_Nord:GW'))*1700|-2>> hectares at 1{,}700\,ha/GW \citep{ISPRA2026}, plus <<UN:ss0_PI_670_h6:PV_agrivoltaico_Sud-isole:GW|1>>\,GW of agrivoltaic, about <<EXPR:V('UN:ss0_PI_670_h6:PV_agrivoltaico_Sud-isole:GW')*2000|-2>> hectares. The system without nuclear requires <<EXPR:V('GR:ss0_PI_none_h6:on:GW')-13.6|1>>\,GW of new onshore wind, about <<EXPR:(V('GR:ss0_PI_none_h6:on:GW')-13.6)*200|-2>> turbines in <<EXPR:(V('GR:ss0_PI_none_h6:on:GW')-13.6)*10|-1>>--<<EXPR:(V('GR:ss0_PI_none_h6:on:GW')-13.6)*20|-1>> parks, and <<EXPR:V('GR:ss0_PI_none_h6:pv:GW')-24.2|1>>\,GW of new solar: <<EXPR:(V('UN:ss0_PI_none_h6:PV_terra_Sud-isole:GW')+V('UN:ss0_PI_none_h6:PV_terra-pipeline_Lazio/Sicilia/Piem./Puglia/Sard.:GW')+V('UN:ss0_PI_none_h6:PV_terra_Centro-Sud:GW')+V('UN:ss0_PI_none_h6:PV_terra_Nord:GW'))|1>>\,GW ground-mounted (about <<EXPR:(V('UN:ss0_PI_none_h6:PV_terra_Sud-isole:GW')+V('UN:ss0_PI_none_h6:PV_terra-pipeline_Lazio/Sicilia/Piem./Puglia/Sard.:GW')+V('UN:ss0_PI_none_h6:PV_terra_Centro-Sud:GW')+V('UN:ss0_PI_none_h6:PV_terra_Nord:GW'))*1700|-2>> hectares), <<EXPR:(V('UN:ss0_PI_none_h6:PV_agrivoltaico_Sud-isole:GW')+V('UN:ss0_PI_none_h6:PV_agrivoltaico_Centro-Sud:GW')+V('UN:ss0_PI_none_h6:PV_agrivoltaico_Nord:GW'))|1>>\,GW of agrivoltaic (about <<EXPR:(V('UN:ss0_PI_none_h6:PV_agrivoltaico_Sud-isole:GW')+V('UN:ss0_PI_none_h6:PV_agrivoltaico_Centro-Sud:GW')+V('UN:ss0_PI_none_h6:PV_agrivoltaico_Nord:GW'))*2000|-3>> hectares at 0.5\,MW per hectare) and <<EXPR:V('GR:ss0_PI_none_h6:pv:GW')-24.2-(V('UN:ss0_PI_none_h6:PV_terra_Sud-isole:GW')+V('UN:ss0_PI_none_h6:PV_terra-pipeline_Lazio/Sicilia/Piem./Puglia/Sard.:GW')+V('UN:ss0_PI_none_h6:PV_terra_Centro-Sud:GW')+V('UN:ss0_PI_none_h6:PV_terra_Nord:GW'))-(V('UN:ss0_PI_none_h6:PV_agrivoltaico_Sud-isole:GW')+V('UN:ss0_PI_none_h6:PV_agrivoltaico_Centro-Sud:GW')+V('UN:ss0_PI_none_h6:PV_agrivoltaico_Nord:GW'))|1>>\,GW of industrial rooftops; ground and agrivoltaic together are about <<EXPR:((V('UN:ss0_PI_none_h6:PV_terra_Sud-isole:GW')+V('UN:ss0_PI_none_h6:PV_terra-pipeline_Lazio/Sicilia/Piem./Puglia/Sard.:GW')+V('UN:ss0_PI_none_h6:PV_terra_Centro-Sud:GW')+V('UN:ss0_PI_none_h6:PV_terra_Nord:GW'))*1700+(V('UN:ss0_PI_none_h6:PV_agrivoltaico_Sud-isole:GW')+V('UN:ss0_PI_none_h6:PV_agrivoltaico_Centro-Sud:GW')+V('UN:ss0_PI_none_h6:PV_agrivoltaico_Nord:GW'))*2000)/1000|-1>> thousand hectares, about <<EXPR:((V('UN:ss0_PI_none_h6:PV_terra_Sud-isole:GW')+V('UN:ss0_PI_none_h6:PV_terra-pipeline_Lazio/Sicilia/Piem./Puglia/Sard.:GW')+V('UN:ss0_PI_none_h6:PV_terra_Centro-Sud:GW')+V('UN:ss0_PI_none_h6:PV_terra_Nord:GW'))*1700+(V('UN:ss0_PI_none_h6:PV_agrivoltaico_Sud-isole:GW')+V('UN:ss0_PI_none_h6:PV_agrivoltaico_Centro-Sud:GW')+V('UN:ss0_PI_none_h6:PV_agrivoltaico_Nord:GW'))*2000)/30.1e6*100|1>>\% of the 30.1 million hectares of the national territory (our estimates; with agrivoltaics farming continues). Individual wind projects meet strong local opposition: in Sardinia a regional moratorium and a law on suitable areas, later struck down by the Constitutional Court \citep{Sardegna2024,Consulta2026}; in the central Apennines 21 projects with 188 turbines contested by local committees \citep{Appennino2025}. Nuclear sites are fewer but each is more contested: in a 2026 Ipsos survey 12\% of Italians would accept a plant within 10\,km and 34\% within 100\,km \citep{Ipsos2026}. A few large consents are administratively simpler than thousands of small ones, but not politically easy. The national repository for radioactive waste, about 95{,}000\,m$^3$ including waste from medicine, research and industry, is needed in any case. It would occupy about 150 hectares, about as much as a single 80\,MW ground-mounted solar park \citep{Deposito2023}: one site, against <<EXPR:(V('GR:ss0_PI_670_h6:on:GW')-13.6)*10|-1>>--<<EXPR:(V('GR:ss0_PI_670_h6:on:GW')-13.6)*20|-1>> wind parks in the central system and <<EXPR:(V('GR:ss0_PI_none_h6:on:GW')-13.6)*10|-1>>--<<EXPR:(V('GR:ss0_PI_none_h6:on:GW')-13.6)*20|-1>> without nuclear. 51 potentially suitable areas have been identified \citep{Deposito2023}, and the only municipality that volunteered withdrew in 2024 \citep{Trino2024}.

\subsection{Supply security and materials}
Fuel is a small and storable share of nuclear cost. A light-water reactor uses about 0.021--0.027\,kg of natural uranium per MWh \citep{WNAecon}; at the long-term price of 96.5\,USD/lb U$_3$O$_8$ \citep{Cameco2026} this is about 5--7\,USD/MWh, and the whole fuel cycle about 11\,USD/MWh, roughly 5\% and 10--12\% of the cost of a new nuclear MWh. Doubling the uranium price raises the cost of a nuclear MWh by about 5\%, whereas fuel dominates the cost of gas-fired power. Production is spread over several countries (Kazakhstan 39\%, Canada 24\%, Namibia 12\%, Australia 8\% in 2024), although Kazakhstan, Uzbekistan and Russia together account for about half \citep{WNAmining}. The critical bottleneck is enrichment, where Rosatom holds 43\% of world capacity and Russia supplied 23\% of the enrichment services delivered to the EU in 2025 \citep{WNAenrich,ESA2025}. Italy has small uranium deposits: according to historical estimates, about 6{,}000\,t of uranium oxide (U$_3$O$_8$) at Val Vedello and about 1{,}000\,t at Novazza \citep{Geopop2026}, together about 5{,}900\,t of uranium, enough for roughly 30\,GW-years at 200\,tU per GW-year. Uranium from seawater is a long-run ceiling on the fuel price: published estimates of 606--1{,}230\,USD/kg \citep{SchneiderGill2014,Schneider2013}, 2.4--5 times today's market price, would add 7--26\,USD/MWh, up to about 23\,EUR/MWh (26 over 1.1188, our calculation), more than the gap of the central case (<<GAP:ss0_PI_none_h3:ss0_PI_670_h3|2>> in PI, <<GAP:ss0_OB_none_h3:ss0_OB_670_h3|2>> in OB, three hours); it is a theoretical ceiling on the price of uranium, not a forecast; the technology has been demonstrated at pilot scale but not commercially. One year of fuel for 1\,GW is about 200\,t of uranium, which is easily stockpiled; EU utilities already hold about 3.6 years of requirements \citep{ESA2025}.

Solar and batteries carry a different dependence. China holds more than 80\% of every stage of the solar supply chain and about 95\% of polysilicon and wafers \citep{IEApv2022}; 98\% of the solar panels imported into the EU in 2024 came from China \citep{Eurostat2025}. China refines about 70\% of lithium and more than 95\% of battery-grade graphite, and produces about 80\% of battery cells \citep{IEAgcmo2025}; it has imposed export controls on graphite, rare earths and battery technologies since 2023, partly suspended until November 2026 \citep{HSFchina2025}. Per MW, onshore wind requires about 10{,}100\,kg of critical minerals, solar 6{,}800\,kg and nuclear 5{,}300\,kg \citep{WNAminerals}; per MWh, given the load factors, nuclear requires several times less. These dependences matter more for the system without nuclear, which in the PI case at six hours needs <<GR:ss0_PI_none_h6:battery:GWh|0>>\,GWh of batteries and <<GR:ss0_PI_none_h6:pv:GW|0>>\,GW of solar, than for the central case, with <<GR:ss0_PI_670_h6:battery:GWh|0>>\,GWh and <<GR:ss0_PI_670_h6:pv:GW|0>>\,GW; the cost of reducing these dependences, for example with components meeting origin requirements, is not quantified by the present model. Nuclear has its own concentrations, in enrichment, fuel fabrication and heavy forgings, which only a handful of firms produce.


\subsection{Industrial base}
Large reactors are typically supplied by a foreign vendor responsible for the design and system integration, with local firms supplying components and balance of plant, as for Dukovany and Barakah. Italy has a supply chain for this role: more than 70 firms with about 2{,}800 direct jobs in nuclear activities in 2022 \citep{TEHA2024}, the refurbishment of Cernavoda unit 1 in Romania led by Ansaldo Nucleare, worth more than 400 million euros \citep{Cernavoda2024}, and the ITER vacuum vessel sectors, the last of which was delivered in October 2026 \citep{ITER2026}. The gaps are nuclear-grade heavy forgings, fuel fabrication \citep{Confindustria2025} and the staff of the safety authority, 79 people in 2026 \citep{ISIN2026}; Ansaldo itself indicates the need for a foreign partner as system integrator \citep{Ansaldo2026}.


\subsection{Political feasibility}
Law 169/2026 delegates the government to issue the implementing decrees within twelve months and does not authorise any plant \citep{Law2026}. The national plan foresees 8--16\,GW of nuclear in 2050, 11--22\% of demand \citep{PNIEC2024}. Compared in shares, our central case has <<RUN:ss0_PI_670_h3:N_GW|1>>\,GW, <<RUN:ss0_PI_670_h3:sh|0>>\% of demand, <<EXPR:('within','above')[min(1,int(V('RUN:ss0_PI_670_h3:N_GW')//16.0001))]>> the range of the plan (8--16\,GW, 11--22\%); on our curve (PI, six hours) the range of 8--16\,GW corresponds to fixN between <<THR6:PI:s:5:16|0>> and <<THR6:PI:s:5:8|0>>\,EUR/kW-yr, that is a construction cost of about <<EXPR:V('CPX:'+str(V('THR6:PI:s:5:16'))+':5:7')|-2>>--<<EXPR:V('CPX:'+str(V('THR6:PI:s:5:8'))+':5:7')|-2>>\,EUR/kW with 7 years of construction; the central cost (8{,}540\,EUR/kW, fixN 670) lies below this range, so building as at Dukovany makes more nuclear economic than the plan foresees. The 2011 referendum repealed the previous nuclear framework with 94\% of votes and 55\% turnout; some constitutional lawyers argue that reintroducing it by law is open to challenge \citep{Cabiddu2026}, while others consider the circumstances changed. Recent polls show a majority of about 55\% in favour of new-generation nuclear \citep{Euromedia2026}. The governing parties and Azione voted for the law; the Democratic Party, the Five Star Movement and AVS voted against \citep{Parties2026}. Regional laws against nuclear plants were struck down by the Constitutional Court in 2010 \citep{Consulta2010}. At EU level, nuclear is in the sustainable finance taxonomy, a challenge by Austria was dismissed in 2025 \citep{Taxonomy2025}, and state aid has been approved for comparable projects \citep{WNN2024}. A programme spanning at least three legislatures needs cross-party support: Germany, Belgium, Spain and Switzerland have changed nuclear policy at least once (Germany in 2010--2011 and 2022; Belgium in 2015, 2022 and 2025; Spain in 2011 and with the extension of Almaraz; Switzerland in 2011 and 2017, with a revision of the ban on new plants approved by Parliament in June 2026 and a referendum expected), although only Belgium has repealed its ban on new plants \citep{WNAcountries}.

\subsection{The role of solar PV}
At the real lifetime cost derived from the main FER~X auction, solar is a large part of every least-cost system, but not of the same kind. In the central case (PI, six hours) the system uses <<EXPR:V('GR:ss0_PI_670_h6:pv:GW')-V('UN:ss0_PI_670_h6:PVex_terra:GW')-V('UN:ss0_PI_670_h6:PVex_tetti:GW')|1>>\,GW of new solar, <<EXPR:V('UN:ss0_PI_670_h6:PV_terra_Sud-isole:GW')+V('UN:ss0_PI_670_h6:PV_terra-pipeline_Lazio/Sicilia/Piem./Puglia/Sard.:GW')+V('UN:ss0_PI_670_h6:PV_terra_Centro-Sud:GW')+V('UN:ss0_PI_670_h6:PV_terra_Nord:GW')|1>> ground-mounted in blocks that are all exhausted and <<UN:ss0_PI_670_h6:PV_agrivoltaico_Sud-isole:GW|1>> agrivoltaic, on top of the 24.2\,GW installed after 2016; without nuclear the PI system uses <<GR:ss0_PI_none_h6:pv:GW|0>>\,GW, of which <<EXPR:V('GR:ss0_PI_none_h6:pv:GW')-24.2|1>> new. Solar is therefore worth building next to nuclear, if it costs what the auctions pay and is not sized for the winter, which nuclear, wind and hydro cover. The plants that set the auction prices are large ground-mounted plants in the best sites. Residential rooftop systems costing 1{,}500--2{,}300\,EUR/kWp \citep{Rooftop2026} correspond to about 100--185\,EUR/MWh at 1{,}066--1{,}300 full-load hours, with 25-year life, 5\% cost of capital and 1.5\% annual O\&M (our calculation), two to four times the cost at which the model builds solar; their private profitability depends on tax credits and avoided network charges, which are transfers from other taxpayers and consumers. Before nuclear is available, every solar MWh displaces mostly gas: the auction price of 57\,EUR/MWh is about half the 2025 average wholesale price of 118\,EUR/MWh \citep{PUN2025}. Ground-mounted solar built now is therefore part of the least-cost path with or without nuclear.

\subsection{Limitations}\label{sec:limits}
The model has one node, no ramp constraints and demand response only through smart charging and interruptible load; it uses four historical weather years, the 2016--2019 load structure and a regime system for 2040, not a transition path. The adequacy reserve, the availability of nuclear (0.97 and 0.87), the shape of heat-pump and vehicle demand and the credits of hydro, pumped hydro and batteries are declared assumptions. Imports have the 2025 average price, exports earn nothing and no help comes from abroad in the critical weeks; the neighbouring systems are not simulated. The sensitivities on CSP, SMRs and gas with capture were computed under the earlier import rule. The import rule cuts both ways. Relative to the capacity-credit rule, net self-sufficiency raises the gap from <<GAP:val_PI_none:c3_PI_670_h6|2>> to <<GAP:ss0_PI_none_h6:ss0_PI_670_h6|2>>\,EUR/MWh (PI, six hours): not counting on other countries favours nuclear, and it is a stated policy choice, not a datum. Conversely, the model assumes that neighbours always absorb Italian exports up to the full interconnection capacity, even when they would also have surplus energy; this favours the system without nuclear, which exports <<UN:ss0_PI_none_h6:IMPORT:TWh|0>>\,TWh a year against <<UN:ss0_PI_670_h6:IMPORT:TWh|0>> in the central system. The nuclear share of demand includes the part of nuclear output that is exported. Hydropower is independent of the simulated weather because its monthly budgets are the same in the four years. The six-hour step underestimates the value of nuclear and absolute cost levels, and one-hour results cannot be computed from three-hour data; worst years exist at six hours only.

Several simplifications favour the system without nuclear: the single node, the absence of stability services, four weather years (with the worst years the gap grows) and technology costs for hydrogen and compressed air that are studies or targets. Others favour nuclear: availability of 0.92 and credit of 0.97 in the reserve, flat bioenergy and geothermal, an evening-heavy but otherwise inflexible demand, full credit for hydro and pumped hydro, and no ramp constraints. Existing solar at operating cost and the free wind block E0 favour renewables in both systems. The OB model is reduced, so OB and PI points differ by model as well as by cost case (without nuclear OB costs <<RUN:ss0_OB_none_h6:cost_new|2>> against <<RUN:ss0_PI_none_h6:cost_new|2>> for the full PI model, six hours); the SMR and wave tests switch those technologies back on, one at a time, on top of the OB base. Our own, unsourced assumptions are the 10\,GW CSP ceiling, the 20-year life and 20\% O\&M of 100-hour storage, the 3-year build of gas with CCS, the 1.287 grid scaling at 7\%, the DAC costs and the 17\% import mark-up. The central Dukovany value, 8{,}540\,EUR/kW, is the gross contract minus the contract fuel (itself disputed between sources) plus owner's costs taken from a US database (INL, 12.57\% of overnight cost; 15--20\% of the EPC cost in a European Commission synthesis \citep{DHaeseleer2013}), not from a Czech estimate; the EPC-only value (7{,}470) gives about <<N6:PI:s:5:604|0>>\,GW. The risk case, 9{,}680\,EUR/kW, is the median of four builds whose conversion to 2025 euros assumes uniform spending (Section~\ref{sec:methods}).

The unit size of the reactors (1.055 to 1.4\,GW) lies within the 10\% reserve margin (7.7\,GW at the three-hour net peak of 76.88\,GW), which covers the loss of five units. Common-mode failures across several units are not modelled.


\section{Conclusions}\label{sec:conclusions}
How much nuclear a zero-emission Italian power system should contain is a function of one quantity, the annualised fixed cost of nuclear fixN, which combines construction cost, cost of capital and construction time (Fig.~\ref{fig:share}), once the system is required not to lean on other countries. For the central case, Dukovany at 8{,}540\,EUR/kW (contract price net of fuel plus owner's costs), seven years and 5\% (fixN 670), the three-hour least-cost system under net self-sufficiency has <<RUN:ss0_PI_670_h3:N_GW|1>>\,GW of nuclear (<<RUN:ss0_PI_670_h3:sh|1>>\% of demand) in PI and <<RUN:ss0_OB_670_h3:N_GW|1>>\,GW (<<RUN:ss0_OB_670_h3:sh|1>>\%) in OB, and the value of nuclear, the cost of the system without it minus the cost with it, is <<GAP:ss0_PI_none_h3:ss0_PI_670_h3|2>> and <<GAP:ss0_OB_none_h3:ss0_OB_670_h3|2>>\,EUR/MWh, lower bounds because the step is still coarse. The import rule matters and is not a technical detail: a plan that counts on net imports or on foreign capacity in the critical weeks looks cheaper because it uses plants that other countries build and pay for. With nuclear, a 10\% net import allowance saves <<EXPR:V('RUN:ss0_PI_670_h6:cost_new')-V('RUN:ss10_PI_670_h6:cost_new')|2>>\,EUR/MWh (PI, six hours) and is a poor target. At six hours nuclear falls below 10\,GW above a fixN of about <<THR6:PI:s:5:10|-1>> (PI) and <<THR6:OB:s:5:10|-1>> (OB), about <<EXPR:V('CPX:'+str(V('THR6:PI:s:5:10'))+':5:7')|-2>> and <<EXPR:V('CPX:'+str(V('THR6:OB:s:5:10'))+':5:7')|-2>>\,EUR/kW with seven years and <<EXPR:V('CPX:'+str(V('THR6:PI:s:5:10'))+':5:15')|-2>> and <<EXPR:V('CPX:'+str(V('THR6:OB:s:5:10'))+':5:15')|-2>> with fifteen. Fifteen years of construction at the central cost leave <<N6:PI:s:5:792.5|0>> and <<N6:OB:s:5:792.5|0>>\,GW, and <<N6:PI:s:5:878|0>> and <<N6:OB:s:5:878|0>>\,GW with the risk cost as well; the worst weather years raise the optimum to about <<RUN:ss0_wx4_PI_670_worst:N_GW|0>>\,GW with a wider gap.

Our recommendation is a nuclear fleet of about <<RUN:ss0_PI_670_h3:N_GW|0>>\,GW, about <<RUN:ss0_PI_670_h3:sh|0>>\% of 2040 demand (PI, three hours; <<RUN:ss0_OB_670_h3:N_GW|0>>\,GW in OB), in a system that is self-sufficient on an annual net basis and in the critical weeks, built at about 8{,}500\,EUR/kW including owner's costs, with about seven years of construction, one series design and engineering completed before construction starts, and financed at a real cost near 5\% (public or guaranteed, as at Dukovany). The condition that changes the recommendation is the fixed cost of nuclear, that is financing and construction time: at a 7\% cost of capital the economic fleet falls to about <<RUN:ss0_PI_7pc_923_h3:N_GW|0>>\,GW (PI) and <<RUN:ss0_OB_7pc_923_h3:N_GW|0>> (OB) at three hours, and fifteen-year construction at the central cost has almost the same effect. Owner's costs remain the least documented input: about <<N6:PI:s:5:604|0>>\,GW without them and about <<N6:PI:s:5:699|0>>\,GW with the upper case (9{,}030\,EUR/kW), at six hours. About <<EXPR:V('RUN:ss0_PI_670_h3:N_GW')/1.055|0>> APR1000 units are needed (about <<EXPR:V('RUN:ss0_PI_670_h3:N_GW')/1.4|0>> APR1400), <<EXPR:V('RUN:ss0_PI_670_h3:N_GW')/1.055/15|1>> a year for 15 years or about <<EXPR:V('RUN:ss0_PI_670_h3:N_GW')/15|1>>\,GW a year against 3.9 in France in 1977--1990, at <<EXPR:V('RUN:ss0_PI_670_h3:N_GW')/1.055/6|0>> to <<EXPR:V('RUN:ss0_PI_670_h3:N_GW')/1.055/4|0>> sites; whether the country can sustain this pace is not determined by the model and has to be verified. Technologies without commercial plants (hydrogen with geological storage, direct air capture, small modular reactors) reduce the value of nuclear only at cost assumptions that no contract supports today.

The model is a regime system for 2040 demand with all fleets already built: it does not describe the situation in 2040. It has one node, no ramp constraints, four weather years and no simulation of neighbouring systems, and its adequacy reserve and nuclear availability are declared assumptions. Where the simplifications favour one option we quantify them in Section~\ref{sec:results}.

\section*{Data and code availability}\label{sec:data}
The model code, input data and results are in the accompanying repository. The v7 model is in \texttt{v7/} (\texttt{code/\allowbreak hl5.py}, \texttt{hl7.py}, \texttt{spec7.py}, \texttt{techs\_v7.json}, \texttt{driver7.py}, \texttt{sweep\_ss.py} for the net self-sufficiency runs; raw runs in \texttt{results.jsonl}), with the reproduction commands in \texttt{v7/README.md}. Results are in \texttt{results/\allowbreak v7.json} and \texttt{results/\allowbreak v7\_sweep.json}; \texttt{v7data.py}, \texttt{tables\_v7.py}, \texttt{make\_figures.py}, \texttt{figband.py} and \texttt{make\_tables.py} generate all figures and tables from the json files, and \texttt{prepare\_data.py} builds the hourly input from Open Power System Data \citep{OPSD2020} and PECD \citep{PECD2021}. The v6 scripts are in \texttt{v6/}.

\section*{Declaration of generative AI and AI-assisted technologies in the writing process}
During the preparation of this work the author used Claude (Anthropic) to search and summarise sources, write the model code, run the calculations and draft the text. The author reviewed and edited the content and takes full responsibility for the content of the publication.

\section*{Declaration of competing interest}
The author declares no known competing financial interests or personal relationships that could have influenced the work reported in this paper.

\section*{Funding}
This research did not receive any specific grant from funding agencies in the public, commercial or not-for-profit sectors.

\FloatBarrier

\appendix
\section*{Supplementary material}
\renewcommand{\thesubsection}{S\arabic{subsection}}
\renewcommand{\thetable}{S\arabic{table}}
\setcounter{table}{0}

\subsection{Monthly analysis}\label{app:monthly}
The monthly analysis of Section~\ref{sec:monthly} uses 2024 monthly output of solar (GSE) and wind (Terna) \citep{GSE2025,Terna2025m}, flat demand of 312.3\,TWh and the real lifetime costs of Table~\ref{tab:inputs}, derived from auction prices with Eq.~\eqref{eq:conv}. The overbuild factor is the ratio of mean monthly output to the output of the worst month; the seasonal deficit is the sum of the shortfalls of monthly output below demand. With solar only the overbuild is 1.95 and the deficit 16.8\% of annual energy (about 52\,TWh); with a 50/50 energy mix of solar and wind they fall to 1.35 and 5\%. Daily storage assumes that 60\% of the energy of the solar share passes through batteries, one cycle per day and 85\% round-trip efficiency, at the real MACSE cost of 11{,}466\,EUR per MWh of capacity per year \citep{MACSE2025}. Table~\ref{tab:monthly} reports the results.

\begin{table}[htbp]
\centering\small
\caption{Monthly analysis: cost per MWh of energy that meets demand in every month (EUR/MWh).}\label{tab:monthly}
\begin{tabular}{p{9.2cm}r}
\toprule
Configuration & EUR/MWh \\
\midrule
Solar only, overbuilt for January, no storage & 82 (96 NZIA) \\
Solar only, overbuilt, daily batteries & 110 (125 NZIA) \\
Solar at mean capacity, seasonal lithium-ion batteries (one cycle per year) & about 2{,}000 \\
Solar, seasonal hydrogen in salt caverns$^a$ & about 390 \\
Solar and wind 50/50, overbuilt, daily batteries & 76--92 (84) \\
\quad with grid cost of 14\,EUR/MWh \citep{Franzmann2026} & 90--106 (98) \\
Best case$^c$, seasonal iron-air storage & about 115 \\
Best case$^c$, overbuilt & about 55 (68 with grid) \\
\bottomrule
\multicolumn{2}{p{11.5cm}}{\footnotesize $^a$ Storage at 37.6\,EUR/kg for purely seasonal use \citep{DCU2026}, 33.3\,kWh/kg, electrolysis 67\%, reconversion 60\%; electrolyser and turbine capital costs excluded, so the value is a lower bound. $^c$ Solar and wind 30\% below the central real costs (29.5 and 44.2\,EUR/MWh), batteries at half the real MACSE cost, iron-air at the 20\,USD/kWh target \citep{Form2026} with 100\% efficiency, no cost for calm periods, heat waves or land.}\\
\end{tabular}
\end{table}

Covering 2024 consumption with overbuilt solar alone would need about 570\,GW, or about 0.97 million hectares at 1{,}700\,ha/GW \citep{ISPRA2026}, 3.2\% of the 30.1 million hectares of the national territory (our calculation); rooftops could host 140--193\,GW \citep{ISPRA2026}, 149--206\,TWh at 1{,}066 h. Existing pumped hydro (7.3\,GW, 53\,GWh in the model) covers about 7 hours at full power \citep{Pumped2025}, roughly one thousandth of the seasonal requirement.

\subsection{Nuclear cost per MWh}\label{app:nuclear}
Table~\ref{tab:lcoe} gives the stand-alone cost per MWh of nuclear at a 90\% load factor, 80-year life, uniform spending during construction and 32\,EUR/MWh for operation, fuel and the back-end fee. In the system model the load factor is endogenous (7{,}473 equivalent hours, 85\%, in the central PI case at three hours), so stand-alone values are lower than the cost per MWh actually produced in the least-cost systems. The formula is $\mathrm{lcoe}=K\phi(r,T_c)\mathrm{CRF}(r,80)/(0.9\times8760)\times1000+32$.

\begin{table}[htbp]
\centering\small
\caption{Stand-alone nuclear cost (EUR/MWh) of the five cases at 90\% load factor, 5\% and 7\%.}\label{tab:lcoe}
\input{tab_lcoe}
\end{table}

Reference construction costs: Shin Kori 3 and 4 (Korea) about 2{,}400\,USD/kW ex post, new Korean units about 3{,}300\,USD/kW \citep{HSF2026,KHNP2025}; Barakah (four APR1400 units) a 20.4 billion USD construction contract (3{,}643\,USD/kW for 5.6\,GW; 4{,}086\,EUR/kW in 2025 euros with the CPI-U factor 1.255 for uniform spending over 2012--2024) and 24.4 billion USD of total financing, which is not used \citep{WNAuae}; the model uses 4{,}670 with owner's costs; IEA assumption for the EU in 2030, 5{,}100\,USD/kW \citep{IEA2024WEO}; Dukovany CZK\,407 billion for two APR1000 units, gross 7{,}813--7{,}850\,EUR/kW (at 24.688 CZK per EUR \citep{EurostatFX}) with fuel for about ten years according to \citet{WNNKHNPplans} and CZK\,15 billion more for the first fuel according to \citet{CT24Dukovany} and \citet{ForbesDukovany}, against the model's central value of 8{,}540, which removes the 382\,EUR/kW of fuel (4.85\,EUR/MWh $\times$ 10 years $\times$ 0.9 $\times$ 8{,}760\,h / 1{,}000; fuel cost of the World Nuclear Association \citep{WNAecon}) and adds owner's costs of 12.57\% of overnight cost \citep{INLowner} \citep{Dukovany2024,WNNDukovanySigned2025,EnerdataDukovany2025}; Flamanville 13.2 billion EUR in 2015 euros excluding interim interest \citep{EDFFlamanville2022} for 1.65\,GW gross \citep{Folz2019}, that is 8{,}000\,EUR/kW in 2015 euros and about 10{,}300 in 2025 euros with the euro-area HICP (factor 1.2875, 128.75 over 100.00; Eurostat dataset prc\_hicp\_aind, updated 6 February 2026, marked discontinued and replaced by prc\_hicp\_ainr, accessed 9 November 2026) \citep{EurostatHICP}; Hinkley Point C about \pounds35 billion in 2015 prices for 3.26\,GW \citep{Hinkley2026}; Vogtle 3 and 4 about 31 billion USD without financing charges \citep{EnergytechVogtle,IEEFAVogtle}, 15{,}656\,EUR/kW in 2025 euros (the model uses 15{,}540), within the 12{,}300--17{,}570\,USD/kW of \citet{Lazard2025}. All five cases are overnight costs in 2025 euros per kW: dollar values of the sources are brought to 2025 with the US consumer price index (CPI-U) and converted at 1 EUR = 1.1188 USD, the Bank of Italy reference rate of 26 October 2026 \citep{BancaItaliaFX}. The risk case, 9{,}680\,EUR/kW with 15 years, is the median of four recent Western builds in 2025 euros (Mochovce 3--4 about 8{,}760, Olkiluoto 3 about 9{,}060, Flamanville 3 about 10{,}300, Vogtle 3--4 15{,}540), calculation and sources in Section~\ref{sec:methods}; its results are interpolated on the curve. The fuel cost of 4.85\,EUR/MWh follows from 4.62\,USD/MWh in September 2021 prices \citep{WNAecon}, the US CPI-U factor 1.174 (274.310 in September 2021, 321.943 in the 2025 average over eleven months \citep{BLSCPI}; annual means before 2016 from the monthly series \citep{FREDCPI}) and 1.1188; the runs used 4.91. Without netting the fuel, Dukovany at 8{,}980\,EUR/kW (fixN 696) gives <<N6:PI:s:5:696|1>>\,GW (PI) and <<N6:OB:s:5:696|1>> (OB) at six hours; without owner's costs, 7{,}470\,EUR/kW (fixN 604) gives <<N6:PI:s:5:604|1>> and <<N6:OB:s:5:604|1>>\,GW at three hours. Interest during construction in our formulation adds 13\% for five years at 5\%, against 17\% in Banca d'Italia's calculation \citep{Bankitalia2025}.

\subsection{Monthly and hourly denominators}\label{app:denom}
The monthly calculation divides cost by total demand. The hourly model divides by demand not covered by hydropower (356.6 of 400\,TWh); on total demand costs are 0.8915 times the hourly figure (356.6 over 400), for example <<EXPR:V('RUN:ss0_PI_670_h3:cost_new')*0.8915|1>> instead of <<RUN:ss0_PI_670_h3:cost_new|2>>\,EUR/MWh for the central PI case at three hours. Hence the monthly figures of Section~\ref{sec:monthly} and the hourly ones are not directly comparable.

\subsection{Input assumptions}\label{app:decl}
The capital share of the real solar cost is 76.4\% of 42.2\,EUR/MWh (the rest is operation, 9.94\,EUR/MWh), and the cost multipliers of the blocks apply to the capital part only. Renewables and storage keep the 5\% rate except in the 7\% variants. The 100-hour storage cost of 1{,}486\,EUR/MWh-yr uses an implicit exchange rate of 1.08\,USD per euro. Assumptions that are our own, not sourced: the 10\,GW ceiling of concentrated solar power, the 20-year life and 20\% O\&M of 100-hour storage, the 3-year build time of gas with CCS, the 1.287 scaling of the grid at 7\%, the DAC costs, and the 17\% mark-up of import prices in the critical weeks. Table~\ref{tab:inputs} lists all inputs.

\subsection{Curve points}
\begin{table}[htbp]
\centering
\caption{The curve under net self-sufficiency at 5\%, six-hour step, with the equivalent construction cost at seven years: system cost (EUR/MWh), nuclear capacity (GW) and share of total demand (\%), PI and OB.}\label{tab:curve5}
\input{tab_curve5}
\end{table}

\begin{table}[htbp]
\centering
\caption{The curve under net self-sufficiency at 7\%, with imports with a capacity credit at 5\% and 7\%, and without imports at 5\% and 7\%, six-hour step.}\label{tab:other}
\input{tab_other}
\end{table}

\begin{table}[htbp]
\centering
\caption{From the previous model to the present one, OB case, six hours, cumulative; cost in EUR/MWh.}\label{tab:bridgeob}
\input{tab_bridge_ob}
\end{table}

\clearpage

\bibliographystyle{plainnat}
\bibliography{refs}
\end{document}
